%% Question bank chapter file - generated by bank_tools.py build (do not edit while the bank is still being built:
%% it is regenerated from the per-test data in Question Bank\data). Every question is one bankq block:
%%   \begin{bankq}[id=..., source={...}, paper=.., qno=.., type=..., subtopic={...}, answer={...}, answer_src=key|solved,
%%                 working={...}, difficulty=1-10, diff_parts={K S I C F}, diff_note={...},
%%                 group=<passage id or empty>, also_in={duplicate copies in other tests}]
%% difficulty = min(10, K+S+I+C+F-1): K concepts 1-3, S steps 1-4, I insight 0-2, C computation 0-2, F format 0-1
%%   ...question LaTeX, options, figures (\bfig / \bfigi -> img/)...
%%   \end{bankq}
%% Passages shared by several questions are bankpassage blocks; their questions carry group=<passage id>.
\documentclass[10pt]{article}
\usepackage[a4paper,margin=1.8cm]{geometry}
\usepackage[T1]{fontenc}
\usepackage{amsmath,amssymb,mathtools}
\usepackage{graphicx}
\usepackage[export]{adjustbox}
\usepackage{xcolor,needspace,keyval}
\usepackage{enumitem}
\usepackage{booktabs,array,tabularx}
\usepackage[hidelinks]{hyperref}
\setlength{\parindent}{0pt}
\setlength{\parskip}{3pt}
\makeatletter
\define@key{bq}{id}{\def\bq@id{#1}}\define@key{bq}{source}{\def\bq@src{#1}}\define@key{bq}{paper}{\def\bq@paper{#1}}
\define@key{bq}{qno}{\def\bq@qno{#1}}\define@key{bq}{type}{\def\bq@type{#1}}\define@key{bq}{subtopic}{\def\bq@sub{#1}}
\define@key{bq}{answer}{\def\bq@ans{#1}}\define@key{bq}{answer_src}{\def\bq@asrc{#1}}\define@key{bq}{working}{\def\bq@work{#1}}
\define@key{bq}{group}{\def\bq@group{#1}}\define@key{bq}{also_in}{\def\bq@also{#1}}
\define@key{bq}{difficulty}{\def\bq@diff{#1}}\define@key{bq}{diff_parts}{\def\bq@dparts{#1}}\define@key{bq}{diff_note}{\def\bq@dnote{#1}}
\newenvironment{bankq}[1][]{\def\bq@work{}\def\bq@also{}\def\bq@group{}\def\bq@diff{?}\def\bq@dparts{}\def\bq@dnote{}\setkeys{bq}{#1}%
  \par\needspace{4\baselineskip}\medskip\noindent{\small\textbf{\bq@id}\hfill\color{gray}[\bq@src, Paper \bq@paper, Q\bq@qno]
  \bq@type\ $\cdot$ \bq@sub\ $\cdot$ \textbf{D\bq@diff}}\par\nopagebreak\relax}
  {\par{\small\color{teal}Answer: \textbf{\bq@ans}\ifx\bq@work\@empty\else\ (solved: \bq@work)\fi
  \ifx\bq@also\@empty\else\par\color{gray}Also in: \bq@also\fi}\par}
\define@key{bp}{id}{\def\bp@id{#1}}\define@key{bp}{source}{\def\bp@src{#1}}\define@key{bp}{questions}{\def\bp@qs{#1}}
\newenvironment{bankpassage}[1][]{\setkeys{bp}{#1}\par\needspace{4\baselineskip}\bigskip\noindent
  {\small\textbf{Passage \bp@id}\hfill\color{gray}[\bp@src; questions \bp@qs]}\par\nopagebreak\relax\itshape}{\par}
\makeatother
\newcommand{\bfig}[2]{\par\smallskip{\centering\includegraphics[width=#1pt,max width=0.62\linewidth,max height=0.32\textheight]{img/#2}\par}\smallskip}
\newcommand{\bfigi}[2]{\includegraphics[width=#1pt,max width=\linewidth,max height=0.32\textheight,valign=c]{img/#2}}
\newcommand{\bmissing}[1]{\textit{[figure missing in the source conversion: \detokenize{#1}]}}
\newenvironment{qoptions}{\begin{itemize}[leftmargin=1.8em,itemsep=1pt,topsep=2pt,label={}]}{\end{itemize}}
\newcommand{\opt}[2]{\item[(#1)] #2}
\providecommand{\fwstart}{}\providecommand{\fwend}{}\providecommand{\qfullsplit}{}
\newcolumntype{L}[1]{>{\raggedright\arraybackslash}p{#1}}
\newcolumntype{C}[1]{>{\centering\arraybackslash}p{#1}}
\newcolumntype{R}[1]{>{\raggedleft\arraybackslash}p{#1}}
\providecommand{\sectionhdr}[1]{}\providecommand{\subsectionhdr}[1]{}


\begin{document}
{\Large\bfseries NLM --- Question Bank}\par\smallskip
{\small Physics chapter P02. Questions from all converted test papers, grouped by source series. Answers: printed key where it matches the paper, otherwise solved (answer\_src=solved, with working). Exact repeats across papers appear once, with also\_in listing the other copies.}\par\medskip\hrule

\section*{Score-I papers}
\begin{bankpassage}[id=SI12-P2-PHY-S1-G3-4, source={Score-I - Test 12 (09 Mar 2022)}, questions={3-4}]
A person stands on a platform-and-pulley system, as shown in Column-I. The masses of the platform and person are $M$ and $m$ respectively. The rope and pulley is massless. The man can pull on the rope so as to produce one of the situations given in following question. In column-I $W_1$ is tension \& $W_2$ is normal reaction between the man and the platform. Match the possible situations.
\begin{center}\small
\begin{tabular}{|p{0.8cm}|p{3.6cm}|p{0.8cm}|p{3.8cm}|p{0.8cm}|p{3.8cm}|}
\hline
 & \centering Column-I & \multicolumn{2}{c|}{Column-II} & \multicolumn{2}{c|}{Column-III}\\
\hline
(I) & \bfigi{66.1}{SI12_p2_phy_s1_para3_4_1.png}\rule{0pt}{35pt} & (i) & $W_2=(2m-M)\dfrac g3$ & (P) & $W_1=(m+M)\dfrac g3$\\
\hline
(II) & \bfigi{77.0}{SI12_p2_phy_s1_para3_4_2.png}\rule{0pt}{45pt} & (ii) & $W_2=(2m-M)\dfrac g2$ & (Q) & $W_1=(m+M)\dfrac g2$\\
\hline
(III) & \bfigi{67.5}{SI12_p2_phy_s1_para3_4_3.png}\rule{0pt}{42pt} & (iii) & $W_2=(m-M)\dfrac g4$ & (R) & $W_1=(m+M)g$\\
\hline
(IV) & \bfigi{52.9}{SI12_p2_phy_s1_para3_4_4.png}\rule{0pt}{33pt} & (iv) & $W_2=(3m-M)\dfrac g4$ & (S) & $W_1=(m+M)\dfrac g4$\\
\hline
\end{tabular}
\end{center}
\end{bankpassage}
\begin{bankq}[id=SI12-P2-PHY-S1-Q03, source={Score-I - Test 12 (09 Mar 2022)}, paper=2, qno=3, type=single, subtopic={Pulley systems}, answer={B}, answer_src=key, difficulty=3, diff_parts={K2 S2 I0 C0 F0}, diff_note={concepts: constraint relations, force balance | steps: FBD for each pulley system; tension relation for platform equilibrium; compare the matching options}, group=SI12-P2-PHY-S1-G3-4]
Platform is in equilibrium then :-
\begin{qoptions}
\opt{A}{(IV) (ii) (Q)}
\opt{B}{(IV) (iv) (R)}
\opt{C}{(IV) (iv) (Q)}
\opt{D}{(III) (ii) (Q)}
\end{qoptions}
\end{bankq}
\begin{bankq}[id=SI12-P2-PHY-S1-Q04, source={Score-I - Test 12 (09 Mar 2022)}, paper=2, qno=4, type=single, subtopic={Pulley systems}, answer={D}, answer_src=key, difficulty=3, diff_parts={K2 S2 I0 C0 F0}, diff_note={concepts: constraint relations, force balance | steps: FBD for each system; acceleration g/2 upward condition; compare the matching options}, group=SI12-P2-PHY-S1-G3-4]
Platform is accelerated by $\dfrac g2$ upward :-
\begin{qoptions}
\opt{A}{(I) (ii) (R)}
\opt{B}{(II) (iv) (Q)}
\opt{C}{(III) (i) (P)}
\opt{D}{(III) (ii) (R)}
\end{qoptions}
\end{bankq}

\section*{TOAS papers}
\begin{bankq}[id=TOAS03-P1-PHY-S1-Q05, source={TOAS - Test 03 (27 Feb 2022)}, paper=1, qno=5, type=single, subtopic={Newton's laws}, answer={B}, answer_src=key, difficulty=3, diff_parts={K2 S2 I0 C0 F0}, diff_note={concepts: pseudo force, motion on an incline | steps: resolve pseudo force and gravity along the incline; relative acceleration; time from s = a t\^{}2/2}]
A block is placed on an inclined plane moving towards right horizontally with an acceleration $a_0=g$. The length of the plane $AC=1$ m. Friction is absent everywhere. The time taken by the block to reach from $C$ to $A$ is ($g=10\ \mathrm{m/s^2}$) :-
\bfig{122.5}{TOAS03_p1_phy_s1_q05_1.png}
\begin{qoptions}
\opt{A}{1.2 s}
\opt{B}{0.74 s}
\opt{C}{2.56 s}
\opt{D}{0.42 s}
\end{qoptions}
\end{bankq}
\begin{bankq}[id=TOAS03-P1-PHY-S2-Q11, source={TOAS - Test 03 (27 Feb 2022)}, paper=1, qno=11, type=multi, subtopic={Friction}, answer={A,B,D}, answer_src=key, difficulty=3, diff_parts={K1 S2 I0 C0 F1}, diff_note={concepts: friction on a wedge and accelerating frame | steps: angle of the incline from the figure; static equilibrium with constant velocity; friction with a = 2; friction with a = 10}]
A wedge is moving rightwards on which a block of mass 10 kg is placed on it. Friction coefficient between the wedge and the block is 0.8. [take $g=10\ \mathrm{m/s^2}$]. Select correct alternative(s) among the following options.
\bfig{163.1}{TOAS03_p1_phy_s2_q11_1.png}
\begin{qoptions}
\opt{A}{If wedge is moving with constant velocity then acceleration of block is zero.}
\opt{B}{If wedge is moving with $\vec a=2\left(\hat i\right)\ \mathrm{m/s^2}$ then friction acting on block is 44N.}
\opt{C}{If wedge is moving with constant velocity then friction acting on block is 64N.}
\opt{D}{If wedge is moving with $\vec a=10\left(\hat i\right)\ \mathrm{m/s^2}$ then friction is 20N, downward on the block along the incline.}
\end{qoptions}
\end{bankq}
\begin{bankq}[id=TOAS03-P1-PHY-S3-Q05, source={TOAS - Test 03 (27 Feb 2022)}, paper=1, qno=5, type=numerical, subtopic={Pulley systems}, answer={3.33}, answer_src=solved, working={Just after the string is cut the spring force is unchanged so the massless pulley stays put; A (m) and C (2m) form an Atwood pair: $a=\frac{2m-m}{3m}g=\frac g3=3.33$ m/s$^2$ ($g=10$), C moves down}, difficulty=4, diff_parts={K1 S2 I1 C0 F1}, diff_note={concepts: pulley system after a string is cut | steps: spring force unchanged just after cutting; massless pulley stays in equilibrium; Atwood pair A and C | I: spring force does not jump}]
In the figure shown, pulley and spring are ideal and strings are light and inextensible. Initially all the bodies are at rest when string connecting $A$ and $B$ is cut. Find initial acceleration of $C$.
\bfig{44.3}{TOAS03_p1_phy_s3_q05_1.png}
\end{bankq}
\begin{bankq}[id=TOAS03-P2-PHY-S1-Q03, source={TOAS - Test 03 (27 Feb 2022)}, paper=2, qno=3, type=multi, subtopic={Friction}, answer={B,C,D}, answer_src=key, difficulty=4, diff_parts={K1 S2 I1 C0 F1}, diff_note={concepts: friction between block and wedge | steps: condition for the block to slip on the wedge; condition for the wedge to move; ranges of mu | I: block and wedge friction conditions}]
A block having mass $m$ is placed over a inclined plane of same mass and inclination $45^\circ$ if coefficient of friction between all surface is same it is $\mu$. Now system set free to move then
\bfig{119.5}{TOAS03_p2_phy_s1_q03_1.png}
\begin{qoptions}
\opt{A}{For $\mu>1$ block over incline will move but wedge don't}
\opt{B}{For $0.2<\mu<1$ block will move but wedge will not}
\opt{C}{For $\mu=0.15$ both will move}
\opt{D}{It is not possible to move wedge without moving block over wedge}
\end{qoptions}
\end{bankq}
\begin{bankq}[id=TOAS03-P2-PHY-S2-Q06, source={TOAS - Test 03 (27 Feb 2022)}, paper=2, qno=6, type=numerical, subtopic={Friction}, answer={7.50}, answer_src=key, difficulty=4, diff_parts={K1 S2 I1 C0 F1}, diff_note={concepts: friction between stacked blocks | steps: limiting friction on the top blocks; common acceleration; force on the bottom blocks; solve F | I: choose the limiting block}]
Four blocks are arranged on a smooth horizontal surface as shown. The masses of the blocks are given (see the figure). The coefficient of static friction between the top and the bottom blocks is $\mu_S$. What is the maximum value of the horizontal force $F$, applied to one of the bottom blocks as shown, that makes all the four blocks move with the same acceleration? [Take : $\mu_S=0.5$, $g=10\ \mathrm{m/s^2}$, $m=1$ kg, $M=2$ kg]
\bfig{201.5}{TOAS03_p2_phy_s2_q06_1.png}
\end{bankq}
\begin{bankq}[id=TOAS03-P2-PHY-S3-Q02, source={TOAS - Test 03 (27 Feb 2022)}, paper=2, qno=2, type=numerical, subtopic={Friction}, answer={9}, answer_src=key, difficulty=3, diff_parts={K1 S2 I0 C0 F1}, diff_note={concepts: sliding of a block on a moving plank | steps: a = mu g = 2 m/s\^{}2; final speed 6 m/s; s = v\^{}2/2a}]
A block of mass 2 kg is gently placed over a massive plank moving horizontally over a smooth surface with velocity 6 m/s. The coefficient of friction between the block and plank is 0.2. The distance travelled by the block till it slides on the plank is ($g=10\ \mathrm{m/s^2}$)
\end{bankq}
\begin{bankq}[id=TOAS04-P1-PHY-S2-Q05, source={TOAS - Test 04 (15 Mar 2022)}, paper=1, qno=5, type=multi, subtopic={Circular dynamics}, answer={A,B,C}, answer_src=key, difficulty=4, diff_parts={K1 S2 I1 C0 F1}, diff_note={concepts: vertical circular motion on an incline | steps: effective gravity g sin(theta); minimum speed at the top; tension at B from energy; tension at C | I: replace g by g sin theta}]
A small sphere of mass $m$ is connected by a string to a nail at $O$ and moves in a circle of radius $r$ on the smooth plane inclined at an angle $\theta$ with the horizontal. If the sphere has a velocity $u$ at the top position $A$. Mark the correct option(s) :-
\bfig{120.2}{TOAS04_p1_phy_s2_q05_1.png}
\begin{qoptions}
\opt{A}{Minimum velocity at $A$ so that string does not get slack instantaneously is $\sqrt{\dfrac35 gr}$.}
\opt{B}{Tension at $B$ is $9mg/5$ if the sphere has required velocity as in option(A).}
\opt{C}{Tension at $C$ is $18\,mg/5$ in situation of option(B).}
\opt{D}{None of these}
\end{qoptions}
\end{bankq}
\begin{bankq}[id=TOAS04-P2-PHY-S1-Q02, source={TOAS - Test 04 (15 Mar 2022)}, paper=2, qno=2, type=multi, subtopic={Circular dynamics}, answer={B,C,D}, answer_src=key, difficulty=4, diff_parts={K2 S2 I0 C0 F1}, diff_note={concepts: circular motion, friction | steps: normal force from radial and gravity components; friction balances the tangential gravity; power of friction}]
A particle of mass $m$ is placed on vertical fixed rough circular track and then it is given velocity $v$ vertically downward at position $A$ on track. If block moves on track with constant speed then
\bfig{147.9}{TOAS04_p2_phy_s1_q02_1.png}
\begin{qoptions}
\opt{A}{Coefficient of friction between block and circular track as function of angle $\theta$ is $\mu=\dfrac{\sin\theta}{\cot\theta+\frac{v^2}{Rg}}$}
\opt{B}{Coefficient of friction between block and circular track as function of angle $\theta$ is $\mu=\dfrac{\sin\theta}{\cos\theta+\frac{v^2}{Rg}}$}
\opt{C}{Instantaneous power due to friction is $(-mg\sin\theta\,v)$}
\opt{D}{Work done from $A$ to $C$ by friction on block will be $(-mgR)$}
\end{qoptions}
\end{bankq}
\begin{bankpassage}[id=TOAS04-P2-PHY-S2-G7-8, source={TOAS - Test 04 (15 Mar 2022)}, questions={7-8}]
A weightless rod of length $\ell$ with a small load of mass $m$ at one of its end is held vertical with its lower end hinged on a horizontal surface. The load touches a wedge of mass $M$ in this position. A slight jerk towards right sets the system in motion (see figure), with rod rotating freely in vertical plane about its lower end. There is no friction.
\bfig{107.5}{TOAS04_p2_phy_s2_para7_8_1.png}
\end{bankpassage}
\begin{bankq}[id=TOAS04-P2-PHY-S2-Q07, source={TOAS - Test 04 (15 Mar 2022)}, paper=2, qno=7, type=single, subtopic={Constraints}, answer={D}, answer_src=key, difficulty=5, diff_parts={K2 S2 I1 C1 F0}, diff_note={concepts: constraint relations, energy conservation | steps: relation between wedge speed and rod angular speed; energy conservation; zero normal force condition; solve M/m | I: constraint between rod and wedge}, group=TOAS04-P2-PHY-S2-G7-8]
For what mass ratio $\dfrac Mm$ will the rod form an angle $\theta=\dfrac\pi3$ with the vertical at the moment the load separates from the wedge ?
\begin{qoptions}
\opt{A}{1}
\opt{B}{2}
\opt{C}{3}
\opt{D}{4}
\end{qoptions}
\end{bankq}
\begin{bankq}[id=TOAS04-P2-PHY-S2-Q08, source={TOAS - Test 04 (15 Mar 2022)}, paper=2, qno=8, type=single, subtopic={Constraints}, answer={D}, answer_src=key, difficulty=3, diff_parts={K2 S2 I0 C0 F0}, diff_note={concepts: constraint relations, energy conservation | steps: constraint at separation; energy conservation; wedge speed}, group=TOAS04-P2-PHY-S2-G7-8]
What is speed of the wedge at that moment ? Neglect friction.
\begin{qoptions}
\opt{A}{$\sqrt{\dfrac{g\ell}2}$}
\opt{B}{$\sqrt{\dfrac{g\ell}4}$}
\opt{C}{$\sqrt{\dfrac{g\ell}6}$}
\opt{D}{$\sqrt{\dfrac{g\ell}8}$}
\end{qoptions}
\end{bankq}
\begin{bankpassage}[id=TOAS04-P2-PHY-S2-G9-10, source={TOAS - Test 04 (15 Mar 2022)}, questions={9-10}]
A long helix made of thin wire is held vertical. The radius and pitch of the helix are $R$ and $\rho$ respectively. A bead begins to slide down the helix. Pitch is defined as displacement of particle along the axis of helix in one rotation. (Given : $\tan\theta=\dfrac{\rho}{2\pi R}$)
\bfig{72.3}{TOAS04_p2_phy_s2_para9_10_1.png}
\end{bankpassage}
\begin{bankq}[id=TOAS04-P2-PHY-S2-Q09, source={TOAS - Test 04 (15 Mar 2022)}, paper=2, qno=9, type=single, subtopic={Circular dynamics}, answer={A}, answer_src=key, difficulty=3, diff_parts={K1 S2 I1 C0 F0}, diff_note={concepts: forces on a bead on a rotating wire | steps: centripetal requirement; gravity components; normal force as the vector sum | I: two perpendicular normal components}, group=TOAS04-P2-PHY-S2-G9-10]
Find the normal force applied by the wire on the bead when the speed of the bead is $v$.
\begin{qoptions}
\opt{A}{$mg\cos\theta\sqrt{1+\left(\dfrac{v^2\cos\theta}{Rg}\right)^2}$}
\opt{B}{$mg\sin\theta\sqrt{1+\left(\dfrac{v^2\cos\theta}{Rg}\right)^2}$}
\opt{C}{$mg\sin\theta\sqrt{1+\left(\dfrac{v^2\sin\theta}{Rg}\right)^2}$}
\opt{D}{$mg\cos\theta\sqrt{1+\left(\dfrac{v^2\sin\theta}{Rg}\right)^2}$}
\end{qoptions}
\end{bankq}
\begin{bankq}[id=TOAS04-P2-PHY-S2-Q10, source={TOAS - Test 04 (15 Mar 2022)}, paper=2, qno=10, type=single, subtopic={Circular dynamics}, answer={B}, answer_src=key, difficulty=4, diff_parts={K2 S2 I1 C0 F0}, diff_note={concepts: friction, circular motion | steps: friction balances the tangential gravity component; normal force from the previous part; mu | I: speed is constant so friction equals gravity along the motion}, group=TOAS04-P2-PHY-S2-G9-10]
Eventually, the bead acquires a constant speed of $v_0$. Find the coefficient of friction between the wire and the bead.
\begin{qoptions}
\opt{A}{$\mu=\dfrac{2\tan\theta}{\sqrt{1+\left(\frac{v^2\sin\theta}{Rg}\right)^2}}$}
\opt{B}{$\mu=\dfrac{\tan\theta}{\sqrt{1+\left(\frac{v^2\cos\theta}{Rg}\right)^2}}$}
\opt{C}{$\mu=\dfrac{\tan\theta}{2\sqrt{1+\left(\frac{v^2\cos\theta}{Rg}\right)^2}}$}
\opt{D}{$\mu=\dfrac{\tan\theta}{2\sqrt{1+\left(\frac{v^2\sin\theta}{Rg}\right)^2}}$}
\end{qoptions}
\end{bankq}
\begin{bankpassage}[id=TOAS04-P2-PHY-S3-G1-2, source={TOAS - Test 04 (15 Mar 2022)}, questions={1-2}]
A tube of mass $M$ hangs from a thread and two balls of mass $m$ slide inside it without friction (see figure). The balls are released simultaneously from the top of the ring and slide down on opposite sides. $\theta$ defines the positions of balls at any time as shown in figure.
\bfig{95.5}{TOAS04_p2_phy_s3_para1_2_1.png}
\end{bankpassage}
\begin{bankq}[id=TOAS04-P2-PHY-S3-Q01, source={TOAS - Test 04 (15 Mar 2022)}, paper=2, qno=1, type=numerical, subtopic={Circular dynamics}, answer={1.50}, answer_src=key, difficulty=5, diff_parts={K2 S2 I1 C0 F1}, diff_note={concepts: ring and beads, normal force | steps: speed of the bead at angle theta; normal force on the ring; vertical condition N cos theta >= Mg; minimise m | I: condition on the bead normal force}, group=TOAS04-P2-PHY-S3-G1-2]
Minimum value of $m$ so that ring will start to rise is $kM$, find value of $k$.
\end{bankq}
\begin{bankq}[id=TOAS04-P2-PHY-S3-Q02, source={TOAS - Test 04 (15 Mar 2022)}, paper=2, qno=2, type=numerical, subtopic={Circular dynamics}, answer={0.66 to 0.67}, answer_src=key, difficulty=4, diff_parts={K1 S2 I1 C0 F1}, diff_note={concepts: tube and bead | steps: speed at theta; normal force; vertical balance for M = 0 gives cos theta | I: tube lifts when the normal force exceeds the weight}, group=TOAS04-P2-PHY-S3-G1-2]
If $M=0$, then the value of angle is $\theta$ at which the tube begins to rise. Find $\cos\theta$.
\end{bankq}
\begin{bankq}[id=TOAS04-P2-PHY-S4-Q07, source={TOAS - Test 04 (15 Mar 2022)}, paper=2, qno=7, type=numerical, subtopic={Circular dynamics}, answer={3}, answer_src=key, difficulty=7, diff_parts={K2 S2 I2 C1 F1}, diff_note={concepts: circular motion, projectile after the string slackens | steps: speed at B from tension zero; energy from A to B; projectile from B must hit C on the line BOC; solve v0\^{}2 | I: collinear B, O, C condition}]
A heavy particle is attached to one end of a light string of length $\ell$ whose other end is fixed at $O$. The particle is projected horizontally with a velocity $v_0$ from its lowest position $A$. When the angular displacement of the string is more than $90^\circ$, the particle leaves the circular path at $B$. The string again becomes taut at $C$ such that $B,O,C$ are colinear. If $v_0^2=\left(\dfrac{4+n\sqrt2}{2}\right)g\ell$, find the value of $n$.
\bfig{76.9}{TOAS04_p2_phy_s4_q07_1.png}
\end{bankq}
\begin{bankq}[id=TOAS05-P1-PHY-S1-Q01, source={TOAS - Test 05 (03 Apr 2022)}, paper=1, qno=1, type=multi, subtopic={Circular dynamics}, answer={A,C}, answer_src=key, difficulty=4, diff_parts={K1 S2 I1 C0 F1}, diff_note={concepts: circular motion with friction and string tension | steps: tension from block B at 0.6 m; limiting friction direction for A; condition for minimum x; condition for maximum x | I: friction direction at the limits}]
Table shown in figure is rotating about its vertical axis with an angular speed $\omega=2$ rad/sec. Blocks $A$ and $B$ have masses 45 kg and 22.5 kg respectively. String connecting them is light and passes through a small hole in the table. Pulley shown in massless. Coefficient of friction at the contacts of $A$ and $B$ with the table is 0.4. dimensions of two blocks are small and distance of $B$ from rotation axis is 0.6 m. If block $A$ remains stationary relative to the table at a distance `$x$' from the rotation axis then (Take : $g=10$ m/s$^2$)
\bfig{97.9}{TOAS05_p1_phy_s1_q01_1.png}
\begin{qoptions}
\opt{A}{Minimum possible value of $x$ is 0.13 m}
\opt{B}{Minimum possible value of $x$ is 0.26 m}
\opt{C}{Maximum possible value of $x$ is 2.37 m}
\opt{D}{Maximum possible value of $x$ is 1.37 m}
\end{qoptions}
\end{bankq}
\begin{bankq}[id=TOAS05-P1-PHY-S1-Q02, source={TOAS - Test 05 (03 Apr 2022)}, paper=1, qno=2, type=multi, subtopic={Circular dynamics}, answer={A,C,D}, answer_src=key, difficulty=4, diff_parts={K2 S2 I0 C0 F1}, diff_note={concepts: conical pendulum with a spring, projectile | steps: vertical balance gives spring force 2Mg; extension and natural length; horizontal equation gives speed; range after the spring breaks}]
A ball of mass $M$ is swing around in a circle in a horizontal plane. On end of spring (spring constant $k$) is fixed and other end is connected to the ball as shown in figure. The ball describes a horizontal circle a distance $h$ above the horizontal floor with constant speed. The stretched spring has a length $\ell$ and makes an angle $60^\circ$ with the vertical. Neglect air resistance.
\bfig{109.9}{TOAS05_p1_phy_s1_q02_1.png}
\begin{qoptions}
\opt{A}{Magnitude of force exerted by spring on block is $2\,Mg$}
\opt{B}{Natural length of spring is $\left(\ell+\dfrac{2Mg}{k}\right)$.}
\opt{C}{Speed of the ball is $\sqrt{\dfrac{3g\ell}{2}}$}
\opt{D}{At some instant of time, if spring breaks, the ball moves a horizontal distance $\sqrt{3h\ell}$, before it hits the floor.}
\end{qoptions}
\end{bankq}
\begin{bankq}[id=TOAS05-P1-PHY-S1-Q03, source={TOAS - Test 05 (03 Apr 2022)}, paper=1, qno=3, type=multi, subtopic={Circular dynamics}, answer={B,C}, answer_src=key, difficulty=3, diff_parts={K1 S2 I0 C0 F1}, diff_note={concepts: loop-the-loop dynamics | steps: speed at B by energy; normal force at B; normal force at C; check separation point}]
A small toy car of mass $m$ slides with negligible friction on a loop track as shown in figure. The toy car starts from rest at a point $H=4R$ above the lowest point of the track. (Acceleration due to gravity is $g$ and radius of circular part is $R$).
\bfig{143.7}{TOAS05_p1_phy_s1_q03_1.png}
\begin{qoptions}
\opt{A}{Speed of toy car at point $B$ is $\sqrt{4gR}$}
\opt{B}{Normal force exerted by track on the toy car at point $B$ is 9 $mg$.}
\opt{C}{Normal force exerted by track on the toy car at position $C$ is 6 $mg$.}
\opt{D}{Toy car looses contact with track at a point between $C$ and $D$.}
\end{qoptions}
\end{bankq}
\begin{bankq}[id=TOAS05-P2-PHY-S1-Q11, source={TOAS - Test 05 (03 Apr 2022)}, paper=2, qno=11, type=single, subtopic={Circular dynamics}, answer={B}, answer_src=key, difficulty=1, diff_parts={K1 S1 I0 C0 F0}, diff_note={concepts: rod pendulum completing a circle | steps: energy needed to reach the top with zero speed; u\^{}2 = gL(2 + sqrt2)}]
A bob of mass $m$ is attached to the lower end of a rigid light rod of length $L$ and other end is fixed to a point $O$ so that the system may rotate in vertical circle about $O$ as shown in the figure. Initially the bob is at point $A$ and rod makes an angle $45^\circ$ from vertical. What minimum speed ($v$) should be imparted to the bob so that it may complete the circular track?
\bfig{70.5}{TOAS05_p2_phy_s1_q11_1.png}
\begin{qoptions}
\opt{A}{$u=\sqrt{gL(2-\sqrt2)}$}
\opt{B}{$u=\sqrt{gL(2+\sqrt2)}$}
\opt{C}{$\sqrt{gL(3+\sqrt2)}$}
\opt{D}{$\sqrt{gL(3-\sqrt2)}$}
\end{qoptions}
\end{bankq}
\begin{bankq}[id=TOAS06-P1-PHY-S2-Q07, source={TOAS - Test 06 (17 Apr 2022)}, paper=1, qno=7, type=multi, subtopic={Circular dynamics}, answer={A,D}, answer_src=key, difficulty=4, diff_parts={K1 S2 I0 C1 F1}, diff_note={concepts: motion on a curved path with speed change | steps: curvature of the parabola at A; normal force from radial equation; friction from the tangential acceleration}]
The 0.8-Mg car travels over the hill having the shape of a parabola. When the car is at point A, it is traveling at 9 m/s and increasing its speed at 3 m/s$^2$. Choose the correct option(s): (Neglect the size of the car)
\bfig{176.8}{TOAS06_p1_phy_s2_q07_1.png}
\begin{qoptions}
\opt{A}{The resultant normal force is 6.86 kN}
\opt{B}{The resultant normal force is 2.73 kN}
\opt{C}{Resultant frictional force that all the wheels of the car exert on the road at this instant is 2.22 kN.}
\opt{D}{Resultant frictional force that all the wheels of the car exert on the road at this instant is 1.18 kN.}
\end{qoptions}
\end{bankq}
\begin{bankq}[id=TOAS06-P1-PHY-S2-Q08, source={TOAS - Test 06 (17 Apr 2022)}, paper=1, qno=8, type=multi, subtopic={Circular dynamics}, answer={B,C}, answer_src=key, difficulty=3, diff_parts={K1 S2 I0 C0 F1}, diff_note={concepts: motion on a smooth ramp | steps: slope angle from the curve at y = 0.2; normal force with centripetal term; tangential acceleration g sin(theta)}]
The 8-kg sack slides down the smooth ramp. If it has a speed of 1.5 m/s when $y=0.2$ m, choose the correct option(s) :
\bfig{181.5}{TOAS06_p1_phy_s2_q08_1.png}
\begin{qoptions}
\opt{A}{Normal reaction the ramp exerts on the sack is 78.44 N}
\opt{B}{Normal reaction the ramp exerts on the sack is 81.84 N}
\opt{C}{The rate of increase in the speed of sack at this instant is 1.96 m/s$^2$}
\opt{D}{The rate of increase in the speed of sack at this instant is 3.92 m/s$^2$}
\end{qoptions}
\end{bankq}
\begin{bankq}[id=TOAS06-P1-PHY-S3-Q04, source={TOAS - Test 06 (17 Apr 2022)}, paper=1, qno=4, type=numerical, subtopic={Circular dynamics}, answer={7.00}, answer_src=key, difficulty=6, diff_parts={K2 S2 I1 C1 F1}, diff_note={concepts: circular motion, energy conservation | steps: string wraps so radius decreases; energy at the top of each turn; minimum speed condition at the last top; initial speed | I: critical point is the top of the smallest circle}]
A string is wrapped over a fixed horizontal cylinder of radius $\dfrac{7}{66}$ m and its other end is connected to a point mass, which hangs in air such that 1.2 m length of string is free. Find minimum velocity (in m/s) given to mass perpendicular to length of cylinder, so that string never gets slack before wrapping over cylinder completely once. (Take $g=9.8$ m/s$^2$)
\bfig{180.5}{TOAS06_p1_phy_s3_q04_1.png}
\end{bankq}
\begin{bankq}[id=TOAS19-P1-PHY-S1-Q06, source={TOAS - Test 19 (05 Feb 2023)}, paper=1, qno=6, type=multi, subtopic={Newton's laws}, answer={A,B,C,D}, answer_src=key, difficulty=4, diff_parts={K1 S2 I1 C0 F1}, diff_note={concepts: equilibrium of a hanging chain | steps: horizontal tension is constant; end tensions from the angles; vertical components give weights; masses and tension at the lowest point}]
ABC is a uniform chain of mass 10 kg with B as lowest point.
\bfig{137.6}{TOAS19_p1_phy_s1_q06_1.png}
\begin{qoptions}
\opt{A}{Tension in chain C $>$ tension at A}
\opt{B}{Tension at B $=48$ N}
\opt{C}{Mass ratio of AB and BC $=9:16$}
\opt{D}{Mass of protion AB is 3.6 kg}
\end{qoptions}
\end{bankq}
\begin{bankq}[id=TOAS19-P1-PHY-S2-Q11, source={TOAS - Test 19 (05 Feb 2023)}, paper=1, qno=11, type=single, subtopic={Friction}, answer={A}, answer_src=key, difficulty=5, diff_parts={K2 S2 I1 C1 F0}, diff_note={concepts: friction at the limit, equilibrium of a mass on cords | steps: slip condition gives the cord angle; geometry of the 5 m cord and the pulley; equilibrium of the hanging mass; angle between the cords}]
A mass of 1kg is supported by a cord 5m long attached to a massless ring free to move on a horizontal fixed rod, as shown in the figure. Coefficient of friction between ring and the rod is 0.75. A second cord is attached to the weight and passes through a pulley at $\dfrac{25}3$m from ring. Block A of mass $m$ is attached to the other end of this cord. Just after released.
\bfig{108.3}{TOAS19_p1_phy_s2_q11_1.png}
\begin{center}\footnotesize\setlength{\tabcolsep}{2pt}
\begin{tabular}{|L{0.8cm}L{3.9cm}|L{0.6cm}L{0.8cm}|}\hline
\multicolumn{2}{|c|}{List-I}&\multicolumn{2}{c|}{List-II}\\\hline
(I)&Tension in the string connected to block A in `N'.&(P)&6\\\hline
(II)&Tension in the 5m cord in `N'&(Q)&8\\\hline
(III)&The value of $\theta$ is $\dfrac\pi k$ radian, possible value of $k$.&(R)&2\\\hline
(IV)&Acceleration of block A in m/s$^2$&(S)&0\\\hline
&&(T)&4\\\hline
&&(U)&12\\\hline
\end{tabular}\end{center}
If the ring is just about to slip
\begin{qoptions}
\opt{A}{I $\to$ P;II $\to$ Q;III $\to$ R;IV $\to$ S}
\opt{B}{I $\to$ Q;II $\to$ P;III $\to$ T;IV $\to$ S}
\opt{C}{I $\to$ T;II $\to$ Q;III $\to$ R;IV $\to$ S}
\opt{D}{I $\to$ U;II $\to$ Q;III $\to$ R;IV $\to$ S}
\end{qoptions}
\end{bankq}
\begin{bankq}[id=TOAS19-P1-PHY-S2-Q12, source={TOAS - Test 19 (05 Feb 2023)}, paper=1, qno=12, type=single, subtopic={Constraints}, answer={D}, answer_src=solved, working={$m=0.6$ kg from the slipping case, so block A is 3 kg. The 1 kg mass starts on the 5 m circle at $(-3,-4)$ from the ring and its allowed direction is exactly along the string to the pulley. $30-T_1=3a$, $T_1-6=a$ give $a=6$ m/s$^2$, $T_1=12$ N; the 5 m cord carries $T_2=8$ N (at rest, radial balance $T_2-8=0$); the cords are at right angles ($k=2$). So I$\to$U, II$\to$Q, III$\to$R, IV$\to$P (the printed key lists C)}, difficulty=6, diff_parts={K2 S2 I2 C1 F0}, diff_note={concepts: constrained motion on a circle, string over a pulley | steps: mass of A from the slipping case; allowed direction of motion; equations for block A and the mass; radial balance}]
A mass of 1kg is supported by a cord 5m long attached to a massless ring free to move on a horizontal fixed rod, as shown in the figure. Coefficient of friction between ring and the rod is 0.75. A second cord is attached to the weight and passes through a pulley at $\dfrac{25}3$m from ring. Block A of mass $m$ is attached to the other end of this cord. Just after released.
\bfig{134.7}{TOAS19_p1_phy_s2_q12_1.png}
\begin{center}\footnotesize\setlength{\tabcolsep}{2pt}
\begin{tabular}{|L{0.8cm}L{3.9cm}|L{0.6cm}L{0.8cm}|}\hline
\multicolumn{2}{|c|}{List-I}&\multicolumn{2}{c|}{List-II}\\\hline
(I)&Tension in the string connected to block A in `N'.&(P)&6\\\hline
(II)&Tension in the 5m cord in `N'&(Q)&8\\\hline
(III)&The value of $\theta$ is $\dfrac\pi k$ radian, possible value of $k$.&(R)&2\\\hline
(IV)&Acceleration of block A in m/s$^2$&(S)&0\\\hline
&&(T)&4\\\hline
&&(U)&12\\\hline
\end{tabular}\end{center}
The value of $m$ was initially chosen such that the ring just began to slide. Now the ring is fixed to the rod and mass of block A is increased to 5m, holding the 1kg mass. Just after the system is released, then
\begin{qoptions}
\opt{A}{I $\to$ P;II $\to$ Q;III $\to$ R;IV $\to$ T}
\opt{B}{I $\to$ Q;II $\to$ P;III $\to$ T;IV $\to$ U}
\opt{C}{I $\to$ T;II $\to$ U;III $\to$ R;IV $\to$ Q}
\opt{D}{I $\to$ U;II $\to$ Q;III $\to$ R;IV $\to$ P}
\end{qoptions}
\end{bankq}
\begin{bankq}[id=TOAS19-P1-PHY-S3-Q03, source={TOAS - Test 19 (05 Feb 2023)}, paper=1, qno=3, type=numerical, subtopic={Circular dynamics}, answer={6.00}, answer_src=key, difficulty=5, diff_parts={K1 S2 I1 C1 F1}, diff_note={concepts: friction on a rotating record on an incline | steps: net force needed toward the axis; worst point on the record; friction over normal reaction; evaluate}]
An old record player of 10 cm radius turns at 10 rad/sec while mounted on a $30^\circ$ incline as in the figure. A particle of mass $m$ can be placed any where on the rotating record. If the least possible coefficient of friction $\mu$ that must exist for no slipping to occur is `$\mu$'. Find $2\sqrt3\mu$.
\bfig{103.2}{TOAS19_p1_phy_s3_q03_1.png}
\end{bankq}

\section*{TOPS papers}
\begin{bankq}[id=TOPS01-P1-PHY-S1-Q01, source={TOPS - Test 01 (31 Jul 2022)}, paper=1, qno=1, type=multi, subtopic={Constraints}, answer={B}, answer_src=key, difficulty=6, diff_parts={K2 S2 I1 C1 F1}, diff_note={concepts: string constraint, Newton's laws | steps: constraint relation between A and B; tension equations for A; tension equations for B; solve for the accelerations | several options to check}]
In the figure shown, mass of $A$ is 5 kg. All pulleys and strings are ideal, then :
\bfig{55.6}{TOPS01_p1_phy_s1_q01_1.png}
\begin{qoptions}
\opt{A}{If acceleration of $A$ is 1 m/s$^2$ upwards, then acceleration of $B$ is 3 m/s$^2$ downwards.}
\opt{B}{If acceleration of $A$ is 3 m/s$^2$ upwards, then acceleration of $B$ is 1 m/s$^2$ downwards.}
\opt{C}{If mass of $B$ is 20 kg, then acceleration of $B$ is $\dfrac{3g}{13}$ downwards.}
\opt{D}{If mass of $B$ is 20 kg, then acceleration of $A$ is $\dfrac{g}{13}$ upwards.}
\end{qoptions}
\end{bankq}
\begin{bankq}[id=TOPS01-P1-PHY-S1-Q02, source={TOPS - Test 01 (31 Jul 2022)}, paper=1, qno=2, type=multi, subtopic={Constraints}, answer={A,C}, answer_src=key, difficulty=6, diff_parts={K2 S2 I1 C1 F1}, diff_note={concepts: string constraint, relative velocity | steps: string length constraints; pulley velocities; block velocity; check each option | four chained constraints}]
In the figure show, the man pulls the strings with velocity $u$ from both ends and the block of mass $M$ always remains horizontal. Then,
\bfig{124.2}{TOPS01_p1_phy_s1_q02_1.png}
\begin{qoptions}
\opt{A}{Velocity of block of mass $M$ is $\dfrac{3u}{4}$.}
\opt{B}{Velocity of pulleys $P_1$ is $2u$.}
\opt{C}{Velocity of pulley $P_2$ is $3u/4$.}
\opt{D}{Velocity of pulley $P_3$ is $u/2$.}
\end{qoptions}
\end{bankq}
\begin{bankq}[id=TOPS01-P1-PHY-S1-Q03, source={TOPS - Test 01 (31 Jul 2022)}, paper=1, qno=3, type=multi, subtopic={Newton's laws}, answer={B,D}, answer_src=key, difficulty=5, diff_parts={K2 S2 I0 C1 F1}, diff_note={concepts: spring force, instantaneous acceleration | steps: equilibrium gives spring forces; cut one spring and keep the other force; net force over mass}]
A block is attached to two identical springs as shown in figure. The extension in upper spring is twice the compression in lower spring. Then :
\bfig{34.5}{TOPS01_p1_phy_s1_q03_1.png}
\begin{qoptions}
\opt{A}{If upper spring is cut, just after that the block will accelerate upward with $\dfrac{2g}{3}$.}
\opt{B}{If upper spring is cut, just after that the block will accelerate downward with $\dfrac{2g}{3}$.}
\opt{C}{If lower spring is cut, just after that the block will accelerate upward with $\dfrac{g}{3}$.}
\opt{D}{If lower spring is cut, just after that the block will accelerate downward with $\dfrac{g}{3}$.}
\end{qoptions}
\end{bankq}
\begin{bankq}[id=TOPS01-P1-PHY-S1-Q04, source={TOPS - Test 01 (31 Jul 2022)}, paper=1, qno=4, type=multi, subtopic={Friction}, answer={D}, answer_src=key, difficulty=6, diff_parts={K2 S2 I1 C1 F1}, diff_note={concepts: pseudo force, friction | steps: friction limit on each block; spring force condition; deduce the natural length case | conceptual with a limit case}]
Two blocks of masses $m_1$ and $m_2$ are connected with massless undeformed spring and placed over a long plank moving with an acceleration `$a$' as shown in figure. The coefficient of friction between the blocks and platform is `$\mu$'.
\bfig{134.7}{TOPS01_p1_phy_s1_q04_1.png}
\begin{qoptions}
\opt{A}{Spring will be stretched if $a>\mu g$}
\opt{B}{Spring will be compressed if $a\le\mu g$}
\opt{C}{Spring will neither be compressed nor be stretched only if $a\le\mu g$}
\opt{D}{Spring will be in its natural length under all conditions}
\end{qoptions}
\end{bankq}
\begin{bankq}[id=TOPS01-P1-PHY-S1-Q05, source={TOPS - Test 01 (31 Jul 2022)}, paper=1, qno=5, type=multi, subtopic={Friction}, answer={A,D}, answer_src=key, difficulty=7, diff_parts={K3 S2 I1 C1 F1}, diff_note={concepts: friction, pseudo force in an elevator, relative motion | steps: effective gravity; check slipping; kinetic friction on each; accelerations | elevator pseudo force}]
A plank of mass $m_1=8$ kg with a bar of mass $m_2=2$ kg placed on its rough surface, lie on a smooth floor of elevator ascending with an acceleration $g/4$. The coefficient of friction is $\mu=1/5$ between $m_1$ and $m_2$. A horizontal force $F=30$ N is applied to the plank. Then, in the reference frame of elevator ($g=10$ m/s$^2$)
\bfig{82.0}{TOPS01_p1_phy_s1_q05_1.png}
\begin{qoptions}
\opt{A}{The acceleration of bar is 2.5 m/s$^2$}
\opt{B}{The acceleration of bar is 3.5 m/s$^2$}
\opt{C}{The acceleration of plank is 5 m/s$^2$}
\opt{D}{The acceleration of plank is $\dfrac{25}{8}$ m/s$^2$.}
\end{qoptions}
\end{bankq}
\begin{bankq}[id=TOPS01-P1-PHY-S1-Q06, source={TOPS - Test 01 (31 Jul 2022)}, paper=1, qno=6, type=multi, subtopic={Friction}, answer={A,C}, answer_src=key, difficulty=8, diff_parts={K3 S3 I1 C1 F1}, diff_note={concepts: friction, string constraint, Newton's laws | steps: constraint between blocks; friction on each surface; equations for both; solve for tension; acceleration | multiple friction surfaces}]
Two blocks with masses $M_1$ and $M_2$ of 10 kg and 20 kg respectively are placed as shown in figure coefficient of friction $=0.2$ between all surfaces, then just after the system is released from rest ($g=10$ m/s$^2$) string and pulley are light. Neglect the rotation of $M$.
\bfig{126.3}{TOPS01_p1_phy_s1_q06_1.png}
\begin{qoptions}
\opt{A}{Tension in string is 306 N nearly.}
\opt{B}{Tension in string is 250 N nearly.}
\opt{C}{Acceleration of $M_2$ block is 4.7 m/s$^2$ nearly.}
\opt{D}{Acceleration of $M_2$ block is 6.5 m/s$^2$ nearly.}
\end{qoptions}
\end{bankq}
\begin{bankq}[id=TOPS01-P1-PHY-S1-Q07, source={TOPS - Test 01 (31 Jul 2022)}, paper=1, qno=7, type=multi, subtopic={Friction}, answer={A,C}, answer_src=key, difficulty=6, diff_parts={K2 S2 I1 C1 F1}, diff_note={concepts: friction, vector resolution | steps: normal reaction from the z component; limiting friction; compare with the applied xy force; direction and contact force | 3D force vector}]
A solid cube of mass 5 kg is placed on a rough horizontal surface, in $xy$--plane as shown. The friction coefficient between the surface and the cube is 0.4. An external force $\vec F=\left(12\hat i+16\hat j+20\hat k\right)$ N is applied on the cube. (use $g=10$ m/s$^2$)
\bfig{97.9}{TOPS01_p1_phy_s1_q07_1.png}
\begin{qoptions}
\opt{A}{The block starts slipping over the surface}
\opt{B}{The friction force on the cube by the surface is 10 N.}
\opt{C}{The friction force acts in $xy$--plane at angle $127^\circ$ with the positive $x$--axis in clockwise direction.}
\opt{D}{The contact force exerted by the surface on the cube is $10\sqrt{10}$ N.}
\end{qoptions}
\end{bankq}
\begin{bankq}[id=TOPS01-P1-PHY-S1-Q08, source={TOPS - Test 01 (31 Jul 2022)}, paper=1, qno=8, type=multi, subtopic={Friction}, answer={B,C,D}, answer_src=key, difficulty=10, diff_parts={K3 S3 I2 C2 F1}, diff_note={concepts: variable friction, integration, equilibrium | steps: friction per element; integrate for total friction; relate to F; tension by integrating a part; evaluate at points | friction varying with position}]
A uniform rod of length `$L$' and mass `$M$' has been placed on a rough horizontal surface. The horizontal force `$F$' applied on the rod at one end such that the rod is just in the state of rest and about to slip. If the coefficient of friction varies according to the relation $\mu=Kx$. (where $x$ is the distance from the opposite end). Where `$K$' is $+$ve constant. Then :
\begin{qoptions}
\opt{A}{Tension at mid point is $F/3$}
\opt{B}{Tension at mid point is $F/4$.}
\opt{C}{Tension at $x=\dfrac L4$ is $\dfrac F{16}$}
\opt{D}{Tension at $x=\dfrac{3L}4$ is $\dfrac{9F}{16}$}
\end{qoptions}
\end{bankq}
\begin{bankpassage}[id=TOPS01-P1-PHY-S2-G9-10, source={TOPS - Test 01 (31 Jul 2022)}, questions={9-10}]
In the arrangement shown in figure the masses $m$ of the bar and $M$ of the wedge, as well as the wedge angle $\alpha$, are known. The masses of the pulley and the thread are negligible. The friction is absent.
\bfig{134.7}{TOPS01_p1_phy_s2_para9_10_1.png}
\end{bankpassage}
\begin{bankq}[id=TOPS01-P1-PHY-S2-Q09, source={TOPS - Test 01 (31 Jul 2022)}, paper=1, qno=9, type=single, subtopic={Constraints}, answer={C}, answer_src=key, difficulty=7, diff_parts={K3 S2 I1 C2 F0}, diff_note={concepts: constraint, pseudo force in the wedge frame, Newton's laws | steps: wedge and bar constraint through the thread; forces on the bar; forces on the wedge; solve the accelerations | moving wedge with a pulley}, group=TOPS01-P1-PHY-S2-G9-10]
Which of the following statement(s) is/are correct ?
\begin{qoptions}
\opt{A}{The acceleration of wedge $M$ is $\dfrac{mg\sin\alpha}{M+m(1-\cos\alpha)}$}
\opt{B}{The acceleration of wedge $M$ is $\dfrac{mg\sin\alpha}{M+2m(1+\cos\alpha)}$}
\opt{C}{The acceleration of bar relative to wedge is $\dfrac{mg\sin\alpha}{M+2m(1-\cos\alpha)}$}
\opt{D}{The acceleration of bar relative to wedge is $\dfrac{mg\sin\alpha}{M+2m(1+\cos\alpha)}$}
\end{qoptions}
\end{bankq}
\begin{bankq}[id=TOPS01-P1-PHY-S2-Q10, source={TOPS - Test 01 (31 Jul 2022)}, paper=1, qno=10, type=single, subtopic={Constraints}, answer={D}, answer_src=key, difficulty=5, diff_parts={K2 S2 I1 C1 F0}, diff_note={concepts: equilibrium, normal reaction | steps: find the tension; wedge equilibrium; friction on the ground; normal on the bar | follows the passage}, group=TOPS01-P1-PHY-S2-G9-10]
If in the previous situation ground were rough to prevent slipping while other surfaces still smooth, then which of the following options is incorrect ?
\begin{qoptions}
\opt{A}{The frictional force between wedge and ground is $mg\sin\alpha$.}
\opt{B}{The tension in the string is $mg\sin\alpha$.}
\opt{C}{Normal force acting on the wedge by the ground is $(M+m)g$}
\opt{D}{Normal force acting on the bar by the wedge is $mg$.}
\end{qoptions}
\end{bankq}
\begin{bankq}[id=TOPS01-P1-PHY-S3-Q01, source={TOPS - Test 01 (31 Jul 2022)}, paper=1, qno=1, type=numerical, subtopic={Constraints}, answer={8.00}, answer_src=key, difficulty=9, diff_parts={K3 S3 I1 C2 F1}, diff_note={concepts: string constraint, pseudo force, Newton's laws | steps: constraint relations; forces on the block; forces on the wedge; solve the system; match the form; a plus b | moving wedge}]
In the arrangement shown in figure, the block of mass $m=2$ kg lies on the wedge of mass $M=8$ kg. The surfaces are smooth and pulley \& strings are massless. If the initial acceleration of wedge is $\dfrac{30\sqrt a}{18+b}$ m/s$^2$, then find $a+b$.
\bfig{127.7}{TOPS01_p1_phy_s3_q01_1.png}
\end{bankq}
\begin{bankq}[id=TOPS01-P1-PHY-S3-Q02, source={TOPS - Test 01 (31 Jul 2022)}, paper=1, qno=2, type=numerical, subtopic={Friction}, answer={162.00}, answer_src=key, difficulty=6, diff_parts={K2 S2 I1 C1 F1}, diff_note={concepts: friction, string tension | steps: limiting friction on A; limiting friction on B including A; tension balance; maximum C}]
The maximum value of mass of block $C$ (in kg) so that neither $A$ nor $B$ moves is (Given that mass of $A$ is 100 kg and that of $B$ is 140 kg. Pulleys are smooth and friction coefficient between $A$ and $B$ and between $B$ and horizontal surface is ($\mu=0.3$) $g=10$ m/s$^2$)
\bfig{120.0}{TOPS01_p1_phy_s3_q02_1.png}
\end{bankq}
\begin{bankq}[id=TOPS01-P1-PHY-S3-Q03, source={TOPS - Test 01 (31 Jul 2022)}, paper=1, qno=3, type=numerical, subtopic={Newton's laws}, answer={4.00}, answer_src=key, difficulty=3, diff_parts={K1 S1 I0 C1 F1}, diff_note={concepts: equilibrium of forces | steps: three unequal non-coplanar forces cannot balance; need four}]
The minimum number of unequal, non-coplanar forces acting on an object so that it remains at rest is :
\end{bankq}
\begin{bankq}[id=TOPS01-P1-PHY-S3-Q04, source={TOPS - Test 01 (31 Jul 2022)}, paper=1, qno=4, type=numerical, subtopic={Friction}, answer={5.00}, answer_src=key, difficulty=5, diff_parts={K2 S1 I1 C1 F1}, diff_note={concepts: friction, minimisation | steps: force at the friction angle; least force formula}]
A heavy body of mass 25 kg is to be dragged along a horizontal plane ($\mu=1/\sqrt3$). If the least force required to move the body is $25x$ newton, find $x$.
\end{bankq}
\begin{bankq}[id=TOPS01-P1-PHY-S3-Q05, source={TOPS - Test 01 (31 Jul 2022)}, paper=1, qno=5, type=numerical, subtopic={Friction}, answer={9.00}, answer_src=key, difficulty=6, diff_parts={K2 S2 I1 C1 F1}, diff_note={concepts: friction, kinematics | steps: compare the force with static friction; deceleration; stopping time; distance travelled | stops before 12 s}]
Initial velocity of the block $\bar v=45\hat i$ while external force on it is $\bar F=25\left(-\hat i\right)$. It coefficient of static and kinetic friction are 0.3 and 0.2 respectively then distance travelled by the block in 12 sec is $25x$ m then $x$ value is ($g=10$ m/s$^2$) :
\bfig{179.6}{TOPS01_p1_phy_s3_q05_1.png}
\end{bankq}
\begin{bankq}[id=TOPS01-P1-PHY-S3-Q06, source={TOPS - Test 01 (31 Jul 2022)}, paper=1, qno=6, type=numerical, subtopic={Newton's laws}, answer={0.00}, answer_src=key, difficulty=8, diff_parts={K3 S2 I1 C2 F1}, diff_note={concepts: friction, pseudo force, spring | steps: wedge and block system; acceleration from the spring force; block equations; spring compression | block on a wedge with a spring}]
A block $A$ of mass $M$ rests on a wedge $B$ of mass $2M$ of inclination $\theta$. There is sufficient friction between $A$ and $B$ so that $A$ does not slip on $B$. If there is no friction between $B$ and ground, the compression in spring is
\bfig{119.6}{TOPS01_p1_phy_s3_q06_1.png}
\end{bankq}
\begin{bankq}[id=TOPS02-P1-PHY-S1-Q03, source={TOPS - Test 02 (28 Aug 2022)}, paper=1, qno=3, type=single, subtopic={Friction}, answer={D}, answer_src=key, difficulty=5, diff_parts={K2 S2 I1 C1 F0}, diff_note={concepts: friction on a rope, tension | steps: acceleration of the whole rope; friction on the half beyond the midpoint; tension equation on that half; solve | distributed friction}]
A rope of length $L$ and $M$ is being pulled on a rough floor by a constant horizontal force $F=Mg$. The force is acting at one end of the rope in the same direction as the length of the rope. The coefficient of kinetic friction between rope and floor is $1/2$. Then, tension at the mid-point of the rope is
\begin{qoptions}
\opt{A}{$\dfrac{Mg}4$}
\opt{B}{$\dfrac{2Mg}5$}
\opt{C}{$\dfrac{Mg}8$}
\opt{D}{$\dfrac{Mg}2$}
\end{qoptions}
\end{bankq}
\begin{bankq}[id=TOPS02-P1-PHY-S1-Q04, source={TOPS - Test 02 (28 Aug 2022)}, paper=1, qno=4, type=single, subtopic={Pulley systems}, answer={A}, answer_src=key, difficulty=6, diff_parts={K2 S3 I1 C1 F0}, diff_note={concepts: string constraint, massless pulley | steps: constraint among A, B and C; massless pulley tension relation; equations for A and B; solve accelerations; check point C}]
All pulleys and threads are ideal. $C$ is a point on the thread as shown. When the system is released from rest, at that instant
\bfig{142.7}{TOPS02_p1_phy_s1_q04_1.png}
\begin{qoptions}
\opt{A}{Acceleration of point $C$ is $5g/2$ upwards}
\opt{B}{Acceleration of block $B$ is $2g$ upwards}
\opt{C}{Acceleration of block $A$ is $g/2$ downwards}
\opt{D}{Acceleration of point $C$ is $g$ downwards}
\end{qoptions}
\end{bankq}
\begin{bankq}[id=TOPS02-P1-PHY-S1-Q05, source={TOPS - Test 02 (28 Aug 2022)}, paper=1, qno=5, type=single, subtopic={Friction}, answer={C}, answer_src=key, difficulty=7, diff_parts={K2 S3 I2 C1 F0}, diff_note={concepts: equilibrium, friction | steps: forces on the large cylinder; normal at the contacts; torque on the small cylinder; friction needed; minimum coefficient | geometry of three cylinders}]
Two identical solid cylinders of radius $r$ are placed touching along their lengths on a horizontal surface. A third smooth cylinder of same material but twice the radius of that of the cylinders is placed lengthwise on them so that the system remains at rest. If all three cylinders have the same length, then minimum value of the coefficient of friction between smaller cylinders and the surface is :
\bfig{82.0}{TOPS02_p1_phy_s1_q05_1.png}
\begin{qoptions}
\opt{A}{$\dfrac1{\sqrt2}$}
\opt{B}{$\dfrac13$}
\opt{C}{$\dfrac1{3\sqrt2}$}
\opt{D}{None of these}
\end{qoptions}
\end{bankq}
\begin{bankq}[id=TOPS02-P1-PHY-S2-Q07, source={TOPS - Test 02 (28 Aug 2022)}, paper=1, qno=7, type=multi, subtopic={Pulley systems}, answer={A,B,D}, answer_src=key, difficulty=9, diff_parts={K3 S3 I2 C1 F1}, diff_note={concepts: string constraint, relative acceleration, Newton's laws | steps: man's acceleration relative to the rope; constraint for the block; tension equations; solve for a1; solve for A0 and T | man climbing a rope}]
All the pulleys are frictionless and massless. String is inextensible and massless. A man of Mass $2M$ is climbing up with acceleration of $a_0$ relative to the point he is holding on the string. A block of mass $M$ is connected with the help of string and pulleys as shown. If tension in the string which man is holding, is $T$. Acceleration of mass $M$ is $a_1$ upwards and acceleration of man with respect to ground is $A_0$ upwards. Then choose the correct options.
\bfig{100.2}{TOPS02_p1_phy_s2_q07_1.png}
\begin{qoptions}
\opt{A}{$a_1=\dfrac{3g+4a_0}9$}
\opt{B}{$A_0=\dfrac{a_0-6g}9$}
\opt{C}{$T=\dfrac{Ma_0+3Mg}9$}
\opt{D}{$T=\dfrac{2Ma_0+6Mg}9$}
\end{qoptions}
\end{bankq}
\begin{bankq}[id=TOPS02-P1-PHY-S2-Q08, source={TOPS - Test 02 (28 Aug 2022)}, paper=1, qno=8, type=multi, subtopic={Newton's laws}, answer={A,B,C}, answer_src=key, difficulty=6, diff_parts={K2 S2 I1 C1 F1}, diff_note={concepts: equilibrium, normal reactions | steps: resolve forces in each case; normals for a V trough; check each statement}]
An iron sphere weighing 10N rests in a V shaped smooth trough whose sides form an angle of $60^\circ$ as shown in the figure. Then the reaction forces are ($g=10$ m/s$^2$)
\bfig{198.1}{TOPS02_p1_phy_s2_q08_1.png}
\begin{qoptions}
\opt{A}{$R_A=10$N \& $R_B=0$ in case (i)}
\opt{B}{$R_A=10$N \& $R_B=10$N in case (ii)}
\opt{C}{$R_A=\dfrac{20}{\sqrt3}$N \& $R_B=\dfrac{10}{\sqrt3}$N in case (iii)}
\opt{D}{$R_A=10$N \& $R_B=10$N in all the three cases}
\end{qoptions}
\end{bankq}
\begin{bankq}[id=TOPS02-P1-PHY-S2-Q09, source={TOPS - Test 02 (28 Aug 2022)}, paper=1, qno=9, type=multi, subtopic={Circular dynamics}, answer={D}, answer_src=key, difficulty=5, diff_parts={K2 S2 I0 C1 F1}, diff_note={concepts: uniform circular motion, acceleration | steps: uniform motion has constant velocity; direction of acceleration changes in circular motion; check the options}]
Which of the following statements is correct?
\begin{qoptions}
\opt{A}{Uniform circular motion may be a uniform motion or a uniformly accelerated motion}
\opt{B}{Non-linear motion cannot be a uniformly accelerated motion as there must be some centripetal acceleration which keeps on changing in direction}
\opt{C}{For an object in circular motion with only two forces acting on it, if one of the forces is constant in magnitude and direction, the other force must be varying in magnitude as well as in direction}
\opt{D}{None of the above}
\end{qoptions}
\end{bankq}
\begin{bankq}[id=TOPS02-P1-PHY-S2-Q12, source={TOPS - Test 02 (28 Aug 2022)}, paper=1, qno=12, type=multi, subtopic={Friction}, answer={A,B,D}, answer_src=key, difficulty=6, diff_parts={K2 S2 I1 C1 F1}, diff_note={concepts: friction, Newton's laws | steps: check whether the blocks slip; limiting friction at each surface; accelerations; friction values | three stacked blocks}]
Three large box are kept stationary over one other as shown in figure. A horizontal force 20 N acts on middle block as shown in figure. Then : ($g=10$ m/s$^2$)
\bfig{152.0}{TOPS02_p1_phy_s2_q12_1.png}
\begin{qoptions}
\opt{A}{Acceleration of lowest block is 4.5 m/s$^2$.}
\opt{B}{Acceleration of upper most block is 2 m/s$^2$.}
\opt{C}{Friction between middle and lower block is 12 N.}
\opt{D}{Friction between middle and upper block is 2 N.}
\end{qoptions}
\end{bankq}
\begin{bankq}[id=TOPS02-P1-PHY-S3-Q04, source={TOPS - Test 02 (28 Aug 2022)}, paper=1, qno=4, type=numerical, subtopic={Pulley systems}, answer={1.20}, answer_src=solved, working={Block pulled right by tension $F$, so $a_b=F/M$. Net horizontal force on the wedge with its two pulleys: slope string segments cancel, leaving $F$ to the left from the horizontal segment and $F\cos37^\circ$ to the right from the free end, so $a_w=0.2F/M$ to the left. Relative acceleration $=1.2F/M$}, difficulty=10, diff_parts={K3 S3 I2 C2 F1}, diff_note={concepts: string constraint, Newton's laws, pulleys on a wedge | steps: forces on the block; forces on the wedge pulleys; slope segments cancel; wedge acceleration; relative acceleration | movable pulleys on a wedge}]
In the arrangement shown, pulleys and string are massless and all surfaces are frictionless. The block and the wedge have same mass $M$. The string is being pulled with a constant force $F$. Relative acceleration between the block and the wedge has a magnitude of $\dfrac{KF}M$. Find $K$.
\bfig{134.5}{TOPS02_p1_phy_s3_q04_1.png}
\end{bankq}
\begin{bankq}[id=TOPS02-P1-PHY-S3-Q06, source={TOPS - Test 02 (28 Aug 2022)}, paper=1, qno=6, type=numerical, subtopic={Circular dynamics}, answer={7.07}, answer_src=solved, working={Friction must supply $mv^2/R$ plus balance $mg\sin\theta$; worst position gives $v^2/R+g\sin\theta=\mu g\cos\theta$, so $v^2=Rg(\mu\cos\theta-\sin\theta)=100(1-\tfrac12)=50$ and $v_0=5\sqrt2\approx7.07$ m/s}, difficulty=10, diff_parts={K3 S3 I2 C2 F1}, diff_note={concepts: circular motion on an incline, friction, components | steps: forces in the plane of the incline; centripetal requirement; worst position; friction limit; solve for the speed | circle in the incline plane}]
A toy car is moving on a circular path of radius 10 m on an inclined plane of inclination $30^\circ$. The coefficient of friction between the tyres and the inclined plane is $\mu=\dfrac2{\sqrt3}$. The maximum constant speed of car, such that car does not slip on the inclined plane is $v_0$ m/s. Assume that width of the car is negligible with respect to radius of circular path. Find the value of $v_0$. ($g=10$ m/s$^2$)
\bfig{108.3}{TOPS02_p1_phy_s3_q06_1.png}
\end{bankq}
\begin{bankq}[id=TOPS02-P2-PHY-S1-Q02, source={TOPS - Test 02 (28 Aug 2022)}, paper=2, qno=2, type=multi, subtopic={Constraints}, answer={A,D}, answer_src=key, difficulty=5, diff_parts={K1 S2 I1 C1 F1}, diff_note={concepts: string constraint | steps: length of string constraint; velocity relation; acceleration relation; angular velocity | ring and block}]
Bead $A$ can slide on a smooth horizontal rod and block $B$ move vertically downward. String connected between them and Pulley are ideal. Block is released gently and at a certain instant (not at the time of releasing) the velocity of ring is $v_1$ and block is $v_2$, At that instant ring is at horizontal distance `$x$' from pulley and vertical distance is `$y$' then mark the correct option(s):
\bfig{127.6}{TOPS02_p2_phy_s1_q02_1.png}
\begin{qoptions}
\opt{A}{$v_2=v_1\cos\theta$}
\opt{B}{$v_1=v_2\cos\theta$}
\opt{C}{$a_2=a_1\cos\theta$, where $a_2$ and $a_1$ represent the acceleration of block and ring respectively at that instant.}
\opt{D}{Angular velocity of the ring with respect to an observer on the pulley is $\dfrac{v_1\sin\theta}{\sqrt{x^2+y^2}}$.}
\end{qoptions}
\end{bankq}
\begin{bankq}[id=TOPS02-P2-PHY-S1-Q05, source={TOPS - Test 02 (28 Aug 2022)}, paper=2, qno=5, type=multi, subtopic={Friction}, answer={A,C}, answer_src=key, difficulty=6, diff_parts={K2 S2 I1 C1 F1}, diff_note={concepts: friction, Newton's laws | steps: direction of friction on each block; force on A or B; numeric cases; check static or kinetic | force cases}]
In the given situation, coefficient of friction between the blocks $A$ and $B$ is $\mu$ and there is no friction between ground and block $B$, acceleration of $A$ is $a_A$ and acceleration of $B$ is $a_B$. A force is acting on either of the blocks, then mark the correct option(s) : (Take $g=10$ m/s$^2$)
\bfig{81.1}{TOPS02_p2_phy_s1_q05_1.png}
\begin{qoptions}
\opt{A}{If acceleration of $A$ and $B$ is in same direction and its magnitude is 2 m/s$^2$ and 3 m/s$^2$ respectively then external force $F$ must act on block $B$.}
\opt{B}{If acceleration of $A$ and $B$ is in same direction and its magnitude is 2 m/s$^2$ and 3 m/s$^2$ respectively, mass of $A$ is 20 kg and mass of $B$ is 50 kg, then friction coefficient between the bodies is 0.4}
\opt{C}{If mass of $A$ is 10 kg and mass of $B$ is 25 kg, $\mu=0.4$ and $F=30$ N is applied on block $A$ then friction between the blocks is static in nature.}
\opt{D}{If mass of $A$ is 10 kg and mass of $B$ is 25 kg, $\mu=0.4$ and $F=250$ N is applied on block $A$ then $a_A=(50/7)$ m/s$^2$.}
\end{qoptions}
\end{bankq}
\begin{bankq}[id=TOPS02-P2-PHY-S2-Q03, source={TOPS - Test 02 (28 Aug 2022)}, paper=2, qno=3, type=numerical, subtopic={Circular dynamics}, answer={1.50}, answer_src=key, difficulty=8, diff_parts={K2 S3 I2 C1 F1}, diff_note={concepts: circular dynamics, pulley systems | steps: tension in the string for B lifting; tension for A lifting; the smaller one governs; tension at the lowest point; solve for v | pendulum with pulleys}]
Initially system is in equilibrium, Now mass $m$ is given velocity $v$ in horizontal direction as shown. Find minimum velocity $v$ in m/s Such that either of block $A$ or $B$ loose contact with ground. given $m_A=m_B=5$kg : ($m=1$ kg, $g=10$ m/s$^2$)
\bfig{90.3}{TOPS02_p2_phy_s2_q03_1.png}
\end{bankq}
\begin{bankq}[id=TOPS02-P2-PHY-S2-Q04, source={TOPS - Test 02 (28 Aug 2022)}, paper=2, qno=4, type=numerical, subtopic={Pulley systems}, answer={1.80}, answer_src=key, difficulty=7, diff_parts={K2 S2 I2 C1 F1}, diff_note={concepts: pulley constraint, Newton's laws | steps: constraint of the movable pulley; equations for each mass; solve; acceleration of the pulley | movable pulley}]
Find the magnitude of acceleration (in m/s$^2$) of pulley of mass 1 kg as shown in figure. Other pulley is massless and string is ideal. (Take : $g=9.8$ m/s$^2$)
\bfig{83.0}{TOPS02_p2_phy_s2_q04_1.png}
\end{bankq}
\begin{bankq}[id=TOPS02-P2-PHY-S2-Q05, source={TOPS - Test 02 (28 Aug 2022)}, paper=2, qno=5, type=numerical, subtopic={Constraints}, answer={4.50}, answer_src=key, difficulty=9, diff_parts={K2 S3 I2 C2 F1}, diff_note={concepts: infinite pulleys, string constraint | steps: constraint for the first pulley; relation between successive pulleys; infinite series; sum; acceleration of the block | infinite series of pulleys}]
In the figure all pulleys are massless and frictionless and strings are light and inextensible. The system consists of infinite number of pulleys and blocks (Blocks are identical and having mass 2 kg). Pulley $P_1$ is moved to left with a constant acceleration of 3 m/s$^2$, find the acceleration of block 1. Assume the strings to remain horizontal.
\bfig{194.5}{TOPS02_p2_phy_s2_q05_1.png}
\end{bankq}
\begin{bankq}[id=TOPS02-P2-PHY-S3-Q02, source={TOPS - Test 02 (28 Aug 2022)}, paper=2, qno=2, type=numerical, subtopic={Circular dynamics}, answer={4.00}, answer_src=key, difficulty=7, diff_parts={K2 S2 I2 C1 F1}, diff_note={concepts: circular dynamics, friction | steps: friction limit; centripetal part at 5 m/s; tangential part from the remaining friction; value | friction circle}]
A racing car driver drives his car on a flat circular track of radius $25/3$ m and a coefficient of friction of 0.5. He drives the car in such a manner that he may attain the maximum possible velocity on the track in a minimum possible time. The magnitude of his tangential acceleration (in SI unit) at an instant when his speed is 5 m/s is
\end{bankq}
\begin{bankq}[id=TOPS02-P2-PHY-S3-Q04, source={TOPS - Test 02 (28 Aug 2022)}, paper=2, qno=4, type=numerical, subtopic={Pulley systems}, answer={5.00}, answer_src=key, difficulty=10, diff_parts={K3 S3 I2 C2 F1}, diff_note={concepts: string constraint, Newton's laws, bead friction | steps: forces on the bead; friction from the string; equations for the blocks; solve the acceleration; evaluate 2p | bead on a string}]
Two blocks of equal mass, 4 kg each, are connected to two ends of a massless string passing over a massless pulley. On one side of the string there is a bead of mass 2kg as shown. When the system is released from rest the bead continues to remain at rest while the two blocks accelerate. If the acceleration of the blocks in SI unit is $p$ then value of $2p$ is. (Take $g=10$m/s$^2$)
\bfig{68.9}{TOPS02_p2_phy_s3_q04_1.png}
\end{bankq}
\begin{bankpassage}[id=TOPS03-P1-PHY-S3-G5-6, source={TOPS - Test 03 (20 Nov 2022)}, questions={5-6}]
Two blocks $A$ and $B$ of masses $m_1=2$kg and $m_2=8$ kg respectively are allowed to move without friction. Block $A$ is on block $B$ and block $B$ slides on smooth fixed inclined plane as shown in the figure. The angle of inclination of inclined plane is $45^\circ$. (Take $g=10$m/s$^2$)
\bfig{145.2}{TOPS03_p1_phy_s3_para5_6_1.png}
\end{bankpassage}
\begin{bankq}[id=TOPS03-P1-PHY-S3-Q05, source={TOPS - Test 03 (20 Nov 2022)}, paper=1, qno=5, type=numerical, subtopic={Constraints}, answer={7.86}, answer_src=solved, working={No friction between A and B, so A has no horizontal acceleration and its vertical acceleration equals that of B along the slope component. Along the slope $(m_2+m_1\sin^2\theta)a=(m_1+m_2)g\sin\theta$ gives $a=\dfrac{10\times10\times0.707}{8+1}\approx7.86$ m/s$^2$}, difficulty=7, diff_parts={K2 S2 I2 C1 F1}, diff_note={concepts: constraint, Newton's laws | steps: constraint for A on B; forces on A; forces on B along the slope; solve acceleration of B | wedge with a block on top}, group=TOPS03-P1-PHY-S3-G5-6]
The acceleration of the block $B$ with respect to the ground in m/s$^2$ is \underline{\hspace{1cm}}
\end{bankq}
\begin{bankq}[id=TOPS03-P1-PHY-S3-Q06, source={TOPS - Test 03 (20 Nov 2022)}, paper=1, qno=6, type=numerical, subtopic={Constraints}, answer={5.56}, answer_src=solved, working={A moves vertically with the vertical component of B: $a_A=a\sin45^\circ=7.86\times0.707\approx5.56$ m/s$^2$}, difficulty=5, diff_parts={K2 S1 I1 C1 F1}, diff_note={concepts: constraint, relative motion | steps: vertical component of B's acceleration; acceleration of A}, group=TOPS03-P1-PHY-S3-G5-6]
The acceleration of the block $A$ with respect to the ground in m/s$^2$ is \underline{\hspace{1cm}}
\end{bankq}
\begin{bankq}[id=TOPS03-P1-PHY-S4-Q08, source={TOPS - Test 03 (20 Nov 2022)}, paper=1, qno=8, type=numerical, subtopic={Friction}, answer={0}, answer_src=key, difficulty=9, diff_parts={K3 S3 I2 C1 F1}, diff_note={concepts: friction, Newton's laws, stacked blocks | steps: check slipping between each pair of blocks; limiting friction values; which surfaces slide; accelerations; result | four blocks with different mu}]
In the figure shown, $\mu_1=0.5$, $\mu_2=0.2$, $\mu_3=0.2$, $\mu_4=0.05$, between the pair of surfaces, from top to bottom. Find the acceleration of third block of mass 30kg.
\bfig{109.7}{TOPS03_p1_phy_s4_q08_1.png}
\end{bankq}
\begin{bankq}[id=TOPS03-P2-PHY-S1-Q06, source={TOPS - Test 03 (20 Nov 2022)}, paper=2, qno=6, type=multi, subtopic={Friction}, answer={A,D}, answer_src=key, difficulty=10, diff_parts={K3 S3 I2 C2 F1}, diff_note={concepts: friction, string constraint, Newton's laws | steps: friction on each block; acceleration of the platform; blocks accelerations; force F; check each | platform with two blocks}]
In the given figure, a long platform of mass $m$ is placed on frictionless surface. Two blocks of masses $4m$ and $m$ ($m=10$kg) are placed on the platform. For both blocks, the coefficient of static friction with platform equal to 0.16 and the coefficient of kinetic friction equal to 0.10. The blocks are connected by a light ideal string through a light pulley (mounted at a movable massless stand ) which is acted upon by an unknown horizontal force $F$. The acceleration of the platform is 2 m/s$^2$. Which of the following options is/are correct? (Take : $g=10$ m/s$^2$)
\bfig{171.1}{TOPS03_p2_phy_s1_q06_1.png}
\begin{qoptions}
\opt{A}{$F$ is 180 N.}
\opt{B}{Acceleration of the $4m$ block is 8 m/s$^2$.}
\opt{C}{Acceleration of the $m$ block is 2 m/s$^2$.}
\opt{D}{The sum of magnitudes of accelerations of the two blocks is 10 m/s$^2$.}
\end{qoptions}
\end{bankq}
\begin{bankq}[id=TOPS03-P2-PHY-S3-Q06, source={TOPS - Test 03 (20 Nov 2022)}, paper=2, qno=6, type=numerical, subtopic={Friction}, answer={4}, answer_src=key, difficulty=10, diff_parts={K3 S3 I2 C2 F1}, diff_note={concepts: pseudo force, friction, relative motion | steps: pseudo force components along and across the rod; normal force; friction; acceleration along the rod; time to travel 4 m | bead on an accelerating rod}]
A bead mass 2 kg is fitted onto a rod with a length of 8 m and can move on it with friction having the coefficient of friction 0.5. At the initial moment the bead is in the middle of the rod. The rod moves translational in a horizontal plane with an acceleration 1 m/s$^2$ in the direction forming an angle $37^\circ$ with the rod in gravity free region. Find the time (in sec) when the bead will leave the rod.
\bfig{124.9}{TOPS03_p2_phy_s3_q06_1.png}
\end{bankq}
\begin{bankpassage}[id=TOPS04-P1-PHY-S3-G1-2, source={TOPS - Test 04 (18 Dec 2022)}, questions={1-2}]
A disc $D$ and a fix vertical cylinder $C$ of radius $R=0.25$ m connected through a horizontal thread $CD$ of length $\ell=2$m lies in a horizontal plane. On imparting an initial velocity of $v=2$ m/s to the disc it starts moving on the path shown in the figure.
\bfig{161.3}{TOPS04_p1_phy_s3_para1_2_1.png}
\end{bankpassage}
\begin{bankq}[id=TOPS04-P1-PHY-S3-Q01, source={TOPS - Test 04 (18 Dec 2022)}, paper=1, qno=1, type=numerical, subtopic={Circular dynamics}, answer={4.00}, answer_src=key, difficulty=7, diff_parts={K2 S2 I2 C1 F1}, diff_note={concepts: string wrapping, circular motion | steps: speed constant, tension; radius shrinks as thread winds; angular speed relation; time summed}, group=TOPS04-P1-PHY-S3-G1-2]
What time (in sec) will the disc take to strike the cylinder ?
\end{bankq}
\begin{bankq}[id=TOPS04-P1-PHY-S3-Q02, source={TOPS - Test 04 (18 Dec 2022)}, paper=1, qno=2, type=numerical, subtopic={Circular dynamics}, answer={8.00}, answer_src=key, difficulty=6, diff_parts={K1 S2 I2 C1 F1}, diff_note={concepts: string wrapping | steps: radius shrinks; arc lengths; total distance}, group=TOPS04-P1-PHY-S3-G1-2]
What is the total distance (in meter) travelled by disc before striking the cylinder ?
\end{bankq}
\begin{bankq}[id=TOPS04-P1-PHY-S4-Q08, source={TOPS - Test 04 (18 Dec 2022)}, paper=1, qno=8, type=numerical, subtopic={Circular dynamics}, answer={2}, answer_src=key, difficulty=10, diff_parts={K3 S3 I2 C2 F1}, diff_note={concepts: energy, circular motion, curvature | steps: speed at C by energy; curvature of the sine curve; centripetal equation; reaction; evaluate}]
A ball of mass $m$ starts moving from rest from the point $O$ along the track $OCB$ (see figure). Find the reaction $R$ (in Newton) on the ball at $C$ if the curve $OCB$ is defined by equation $y=h\sin\left(\dfrac{\pi x}\ell\right)$ where $h=\dfrac\ell3$. (Use : $mg=\dfrac{18}{9+2\pi^2}$ Newton)
\bfig{104.3}{TOPS04_p1_phy_s4_q08_1.png}
\end{bankq}
\begin{bankq}[id=TOPS04-P2-PHY-S1-Q01, source={TOPS - Test 04 (18 Dec 2022)}, paper=2, qno=1, type=single, subtopic={Friction}, answer={D}, answer_src=key, difficulty=5, diff_parts={K2 S2 I1 C1 F0}, diff_note={concepts: friction, Newton's laws | steps: string tension and block D; condition for A; condition for B and C; maximum}]
Three blocks $A$, $B$ and $C$ of equal mass $m$ are placed one over other on a frictionless surface (table) as shown in the figure. Coefficient of friction between any blocks $A$, $B$ and $C$ is $\mu$. The maximum value of mass of block $M_D$ so that the block $A$, $B$ and $C$ move without slipping over each other is :-
\bfig{108.3}{TOPS04_p2_phy_s1_q01_1.png}
\begin{qoptions}
\opt{A}{$\dfrac{3m\mu}{\mu+1}$}
\opt{B}{$\dfrac{3m(1-\mu)}\mu$}
\opt{C}{$\dfrac{3m(1+\mu)}\mu$}
\opt{D}{$\dfrac{3m\mu}{(1-\mu)}$}
\end{qoptions}
\end{bankq}
\begin{bankq}[id=TOPS04-P2-PHY-S1-Q14, source={TOPS - Test 04 (18 Dec 2022)}, paper=2, qno=14, type=single, subtopic={Newton's laws}, answer={A}, answer_src=key, difficulty=6, diff_parts={K2 S2 I2 C1 F0}, diff_note={concepts: chain on a sphere, Newton's law | steps: length on the sphere; hanging part; net force; acceleration}]
A uniform flexible chain of length 1.5 m is placed on a smooth spherical surface of radius $R=\left(\dfrac2\pi\right)$ m with its upper end fixed at the top of the sphere while the other end of the chain is hanging freely as shown in figure. Find the acceleration of each element of the chain when its upper end is released.
\bfig{45.6}{TOPS04_p2_phy_s1_q14_1.png}
\begin{qoptions}
\opt{A}{$\left(1+\dfrac4\pi\right)\dfrac g3$}
\opt{B}{$\left(1+\dfrac2\pi\right)\dfrac g6$}
\opt{C}{$\left(1-\dfrac2\pi\right)\dfrac g3$}
\opt{D}{$\left(1+\dfrac3\pi\right)\dfrac g2$}
\end{qoptions}
\end{bankq}
\begin{bankq}[id=TOPS04-P2-PHY-S1-Q17, source={TOPS - Test 04 (18 Dec 2022)}, paper=2, qno=17, type=single, subtopic={Circular dynamics}, answer={A}, answer_src=key, difficulty=6, diff_parts={K2 S2 I2 C1 F0}, diff_note={concepts: friction on a curve, optimisation | steps: speed limit with variable friction; maximise in r; position}]
A cyclist moves along a horizontal circular road of radius $R$, the friction coefficient being dependent only on the distance $r$ from the centre of the circular road as $\mu=\mu_0\left(1-\dfrac rR\right)$, where $\mu_0$ is a constant. What is the value of $r$ for which the cyclist can move with maximum speed without skidding ?
\begin{qoptions}
\opt{A}{$\dfrac R2$}
\opt{B}{$\dfrac R3$}
\opt{C}{$\dfrac R4$}
\opt{D}{$\dfrac R5$}
\end{qoptions}
\end{bankq}
\begin{bankq}[id=TOPS04-P2-PHY-S2-Q03, source={TOPS - Test 04 (18 Dec 2022)}, paper=2, qno=3, type=numerical, subtopic={Circular dynamics}, answer={1}, answer_src=solved, working={On a smooth cone the speed satisfies $v^2=gh=9.8\times0.098=0.96$, so $v\approx0.98\approx1$ m/s}, difficulty=6, diff_parts={K2 S2 I1 C1 F1}, diff_note={concepts: conical surface, circular motion | steps: forces on the particle; radius height relation; speed}]
A particle describes a horizontal circle on a smooth inner surface of an inverted cone. The height of the plane of the circle above the vertex is 9.8 cm. Find the speed (in m/s) of the particle. (Roundoff the answer to closest integer.) ($g=9.8$ m/s$^2$)
\end{bankq}
\begin{bankq}[id=TOPS04-P2-PHY-S2-Q04, source={TOPS - Test 04 (18 Dec 2022)}, paper=2, qno=4, type=numerical, subtopic={Circular dynamics}, answer={20}, answer_src=solved, working={Sliding limit: $a_t^2+(v^2/r)^2=(\mu g)^2=4$ gives $v^2/r=\sqrt2$, $v^2=40\sqrt2$. Then $s=v^2/(2a_t)=20$ m}, difficulty=7, diff_parts={K2 S2 I2 C1 F1}, diff_note={concepts: friction on a curve, kinematics | steps: friction limit; tangential and centripetal parts; speed at sliding; distance}]
A car starts moving with constant tangential acceleration $a_t=\sqrt2$ m s$^{-2}$ along a horizontal circular road of radius $r=40$ m. The coefficient of sliding friction between the wheels of the car and the road is $\mu=0.20$. Find the distance (in m) moved by the car without sliding. (Round off the answer to closest integer.) ($g=10$ m/s$^2$)
\end{bankq}
\begin{bankq}[id=TOPS04-P2-PHY-S2-Q05, source={TOPS - Test 04 (18 Dec 2022)}, paper=2, qno=5, type=numerical, subtopic={Circular dynamics}, answer={18.00}, answer_src=key, difficulty=9, diff_parts={K2 S3 I2 C2 F1}, diff_note={concepts: string through a ring, circular motion | steps: tension equal in both parts; geometry; vertical balance; centripetal equation; speed form; a and b}]
A single string $ACB$ of length 2 m passes through a ring at $C$ that is attached to a sphere which revolves at a constant speed $v$ in the horizontal circle as shown in Figure. Given $\theta_1=60^\circ$, $\theta_2=30^\circ$, $g=10$ ms$^{-2}$, the string is light and inextensible, and tension is same in both parts of the string, If the speed $v$ (in m/s) is $\sqrt{\dfrac{10\sqrt a}{\sqrt b+1}}$. Then find $a^2+b^2$.
\bfig{110.4}{TOPS04_p2_phy_s2_q05_1.png}
\end{bankq}
\begin{bankq}[id=TOPS04-P2-PHY-S2-Q09, source={TOPS - Test 04 (18 Dec 2022)}, paper=2, qno=9, type=numerical, subtopic={Circular dynamics}, answer={2}, answer_src=solved, working={Energy: $v^2=gR(\phi-\sin\phi)$. Leaving the pulley needs $N=0$: $g\sin\phi=v^2/R$, so $2\sin\phi=\phi$ and $\phi/\sin\phi=2$}, difficulty=8, diff_parts={K2 S2 I2 C2 F1}, diff_note={concepts: circular motion, energy | steps: energy of the two masses; condition of leaving; combine; ratio}]
Figure shows a smooth cylindrical pulley of radius $R$ with centre at origin of co-ordinates. An ideal thread is thrown over it on the two parts of ideal thread two identical masses are tied initially at rest with co-ordinates $(R,0)$ and $(-R,-R)$ respectively. If mass at $x$-axis is given a slight upward jerk, it leaves contact with pulley at $(R\cos\phi\ R\sin\phi)$. Then find $\phi/\sin\phi$.
\bfig{111.1}{TOPS04_p2_phy_s2_q09_1.png}
\end{bankq}
\begin{bankq}[id=TOPS04-P2-PHY-S2-Q10, source={TOPS - Test 04 (18 Dec 2022)}, paper=2, qno=10, type=numerical, subtopic={Circular dynamics}, answer={1.00}, answer_src=key, difficulty=10, diff_parts={K3 S3 I2 C2 F1}, diff_note={concepts: rotating cord, tension, elasticity | steps: stretched length; tension at the verge; speed relation; substitute; value}]
An elastic cord having an upstretched length $\ell$, force constant $k$ and mass per unit length $m_0$ is stretched around the drum of radius $R(2\pi R>1)$. Determine the speed of the cord, due to rotation of the drum, which will allow the cord to loosen its contact with the drum. Express value of $\dfrac{v^2}{11}$ (in m$^2$/s$^2$). (for the given data $m_0=40$ g/cm, $k=100$ N/m, $\pi=\dfrac{22}7$, $\ell=400$ cm, $R=70$ cm)
\bfig{76.5}{TOPS04_p2_phy_s2_q10_1.png}
\end{bankq}
\begin{bankq}[id=TOPS05-P1-PHY-S1-Q03, source={TOPS - Test 05 (14 Jan 2023)}, paper=1, qno=3, type=multi, subtopic={Constraints}, answer={A,C}, answer_src=key, difficulty=7, diff_parts={K2 S2 I2 C1 F1}, diff_note={concepts: constraint, wedge motion | steps: constraint between block and wedge; forces on each; accelerations}]
In the given figure, all surface are frictionless and strings pulleys are massless, then
\bfig{124.2}{TOPS05_p1_phy_s1_q03_1.png}
\begin{qoptions}
\opt{A}{acceleration of block $B$ is $g/2$.}
\opt{B}{acceleration of block $B$ is zero.}
\opt{C}{acceleration of block wedge is zero.}
\opt{D}{acceleration of block wedge is $g/2$.}
\end{qoptions}
\end{bankq}
\begin{bankq}[id=TOPS05-P1-PHY-S3-Q02, source={TOPS - Test 05 (14 Jan 2023)}, paper=1, qno=2, type=numerical, subtopic={Pulley systems}, answer={0.305}, answer_src=key, difficulty=8, diff_parts={K3 S2 I2 C1 F1}, diff_note={concepts: pulleys, friction, time varying force | steps: motion condition from friction; force at the moment; equation of each block; acceleration}]
Two identical blocks each having a mass of 20 kg are connected to each other by a light inextensible string as shown and are placed over a rough surface. Pulleys are connected to the blocks. Find acceleration $a_1$ (in m/s$^2$) of the blocks after one second after the application of the time varying force of $40\,t$ N, where $t$ is in second.
\bfig{190.3}{TOPS05_p1_phy_s3_q02_1.png}
\end{bankq}
\begin{bankq}[id=TOPS05-P2-PHY-S2-Q10, source={TOPS - Test 05 (14 Jan 2023)}, paper=2, qno=10, type=multi, subtopic={Constraints}, answer={A,B,C}, answer_src=key, difficulty=9, diff_parts={K3 S3 I2 C1 F1}, diff_note={concepts: pseudo force, pulley, friction | steps: acceleration of the string; forces on the block in the car frame; tension; friction; acceleration}]
A flat car is given an acceleration $a_0$ and started from rest as shown in the figure. Neglect friction between floor and car wheels and also the mass of the pulley. Mass of block $A=50$ kg. Then :-
\bfig{123.9}{TOPS05_p2_phy_s2_q10_1.png}
\begin{qoptions}
\opt{A}{Tension in the string connected to block $A$ is 350 N.}
\opt{B}{Frictional force acting on $A$ is 150 N.}
\opt{C}{Acceleration of block $A$ is 4 m/s$^2$.}
\opt{D}{Acceleration of block $A$ is 1 m/s$^2$.}
\end{qoptions}
\end{bankq}

\section*{TNAS papers}
\begin{bankq}[id=TNAS01-P2-PHY-S2-Q11, source={TNAS - Test 01 (05 Feb 2023)}, paper=2, qno=11, type=single, subtopic={Friction}, answer={A}, answer_src=key, difficulty=5, diff_parts={K1 S2 I2 C1 F0}, diff_note={concepts: ring on a rod with friction | steps: geometry of the cords; equilibrium; tensions; angle}]
A mass of $1$kg is supported by a cord $5$m long attached to a massless ring free to move on a horizontal fixed rod, as shown in the figure. Coefficient of friction between ring and the rod is $0.75$. A second cord is attached to the weight and passes through a pulley at $\dfrac{25}{3}$m from ring. Block A of mass $m$ is attached to the other end of the cord. Just after released.
\bfig{124.2}{TNAS01_p2_phy_s2_q11_1.png}
\begin{center}
{\footnotesize\setlength{\tabcolsep}{2pt}\begin{tabular}{|c|p{4.2cm}|c|p{1.4cm}|}\hline
 & \multicolumn{1}{c|}{List-I} & & List-II\\\hline
(I) & Tension in the string connected to block A in `N'. & (P) & $6$\\\hline
(II) & Tension in the $5$m cord in `N' & (Q) & $8$\\\hline
(III) & The value of $\theta$ is $\dfrac\pi k$ radian, possible value of $k$. & (R) & $2$\\\hline
(IV) & Acceleration of block A in m/s$^2$ & (S) & $0$\\\hline
 & & (T) & $4$\\\hline
 & & (U) & $12$\\\hline
\end{tabular}}
\end{center}
If the ring is just about to slip
\begin{qoptions}
\opt{A}{I $\to$ P;II $\to$ Q;III $\to$ R;IV $\to$ S}
\opt{B}{I $\to$ Q;II $\to$ P;III $\to$ T;IV $\to$ S}
\opt{C}{I $\to$ T;II $\to$ Q;III $\to$ R;IV $\to$ S}
\opt{D}{I $\to$ U;II $\to$ Q;III $\to$ R;IV $\to$ S}
\end{qoptions}
\end{bankq}
\begin{bankq}[id=TNAS01-P2-PHY-S2-Q12, source={TNAS - Test 01 (05 Feb 2023)}, paper=2, qno=12, type=single, subtopic={Pulley systems}, answer={C}, answer_src=key, difficulty=5, diff_parts={K1 S2 I2 C1 F0}, diff_note={concepts: pulley system with a fixed ring | steps: new masses; equations; tensions; acceleration}]
A mass of $1$kg is supported by a cord $5$m long attached to a massless ring free to move on a horizontal fixed rod, as shown in the figure. Coefficient of friction between ring and the rod is $0.75$. A second cord is attached to the weight and passes through a pulley at $\dfrac{25}{3}$m from ring. Block A of mass $m$ is attached to the other end of this cord. Just after released.
\bfig{134.7}{TNAS01_p2_phy_s2_q12_1.png}
\begin{center}
{\footnotesize\setlength{\tabcolsep}{2pt}\begin{tabular}{|c|p{4.2cm}|c|p{1.4cm}|}\hline
 & \multicolumn{1}{c|}{List-I} & & List-II\\\hline
(I) & Tension in the string connected to block A in `N'. & (P) & $6$\\\hline
(II) & Tension in the $5$m cord in `N' & (Q) & $8$\\\hline
(III) & The value of $\theta$ is $\dfrac\pi k$ radian, possible value of $k$. & (R) & $2$\\\hline
(IV) & Acceleration of block A in m/s$^2$ & (S) & $0$\\\hline
 & & (T) & $4$\\\hline
 & & (U) & $12$\\\hline
\end{tabular}}
\end{center}
The value of $m$ was initially chosen such that the ring just began to slide. Now the ring is fixed to the rod and mass of block A is increased to $5m$, holding the $1$kg mass. Just after the system is released, then
\begin{qoptions}
\opt{A}{I $\to$ P;II $\to$ Q;III $\to$ R;IV $\to$ T}
\opt{B}{I $\to$ Q;II $\to$ P;III $\to$ T;IV $\to$ U}
\opt{C}{I $\to$ T;II $\to$ U;III $\to$ R;IV $\to$ Q}
\opt{D}{I $\to$ U;II $\to$ Q;III $\to$ R;IV $\to$ P}
\end{qoptions}
\end{bankq}
\begin{bankq}[id=TNAS03-P1-PHY-S1-Q14, source={TNAS - Test 03 (28 May 2023)}, paper=1, qno=14, type=single, subtopic={Friction}, answer={B}, answer_src=key, difficulty=5, diff_parts={K1 S2 I2 C1 F0}, diff_note={concepts: equilibrium of leaning boards | steps: torque about the top; friction at the bottom; limiting angle}]
Two uniform boards, tied together with the help of string are balanced on a surface as shown. The coefficient of static friction between boards and surface is $0.5$. The minimum value for $\theta$, for which this type of arrangement is possible is :
\bfig{79.6}{TNAS03_p1_phy_s1_q14_1.png}
\begin{qoptions}
\opt{A}{$30^\circ$}
\opt{B}{$45^\circ$}
\opt{C}{$37^\circ$}
\opt{D}{$53^\circ$}
\end{qoptions}
\end{bankq}

\section*{Star Review (Other) papers}
\begin{bankq}[id=STAR01-P1-PHY-S1-Q04, source={Star Review (Other) - Test 01 (12 Sep 2022)}, paper=1, qno=4, type=multi, subtopic={Circular dynamics}, answer={A,B,C,D}, answer_src=key, difficulty=9, diff_parts={K2 S3 I2 C2 F1}, diff_note={concepts: vertical loops, minimum speed and normal force | steps: condition at the top; energy conservation to the top; case two with the different geometry; normal force at B; compare each statement}]
A small block of mass 3 kg is projected into a vertical return loop at A with velocity $v$. The block travels without friction along a circle of radius 2m, and deposited on a horizontal surface at C for each of two loops shown :
\bfig{163.9}{STAR01_p1_phy_s1_q04_1.png}
\begin{qoptions}
\opt{A}{In case-I, the smallest velocity for which the block will reach the horizontal surface at C is given by 10 m/s.}
\opt{B}{In case-II, the smallest velocity for which the block will reach the horizontal surface at C is given by $\sqrt{80}$m/s.}
\opt{C}{In case-I force exerted by the loop on the block as it passes point B is given by 90 N.}
\opt{D}{In case-II force exerted by the loop on the block as it passes point B is given by 60 N.}
\end{qoptions}
\end{bankq}
\begin{bankq}[id=STAR01-P1-PHY-S1-Q05, source={Star Review (Other) - Test 01 (12 Sep 2022)}, paper=1, qno=5, type=multi, subtopic={Pulley systems}, answer={A,B,D}, answer_src=key, difficulty=6, diff_parts={K1 S2 I2 C1 F1}, diff_note={concepts: accelerations of two masses on a pulley pulled by a force | steps: tension from the pulley force; acceleration of each mass; check whether the lighter lifts; compare statements}]
Two blocks of masses $m_1=2$ kg and $m_2=4$ kg hang over a massless pulley as shown in the figure. A force $F_0=100$ N acting at the axis of the pulley accelerates the system upwards. Then :
\bfig{81.6}{STAR01_p1_phy_s1_q05_1.png}
\begin{qoptions}
\opt{A}{Acceleration of 2 kg mass is 15 m/s$^2$}
\opt{B}{Acceleration of 4 kg mass is 2.5 m/s$^2$}
\opt{C}{Acceleration of both the masses is same}
\opt{D}{Acceleration of both the masses is upward.}
\end{qoptions}
\end{bankq}
\begin{bankq}[id=STAR01-P1-PHY-S1-Q06, source={Star Review (Other) - Test 01 (12 Sep 2022)}, paper=1, qno=6, type=multi, subtopic={Circular dynamics}, answer={A,C,D}, answer_src=key, difficulty=9, diff_parts={K2 S3 I2 C2 F1}, diff_note={concepts: motion of a body on an incline tied by a thread, friction | steps: work energy at an angle; speed at the angle; radial equation for tension; maximise the tension; maximum speed condition; check statements}]
One end of a thread of length `$\ell$' is attached to an inclined plane of elevation of $\alpha$. A point like body of mass $m=1$ kg is attached to the other end of the thread as shown in the figure. The body is released without initial speed, such that the thread is tight and horizontal. The coefficient of friction between the slope and body is $\mu$.
\bfig{124.2}{STAR01_p1_phy_s1_q06_1.png}
\begin{qoptions}
\opt{A}{The tension in the thread, when the angle between the thread and the horizontal is $\phi$, is $T=mg\left(3\sin\phi\sin\alpha-2\mu\phi\cos\alpha\right)$.}
\opt{B}{The tension in the string will be maximum when $\phi=90^\circ$.}
\opt{C}{The tension in the string will be maximum when $\phi=\cos^{-1}\left(\dfrac{2\mu}{3\tan\alpha}\right)$.}
\opt{D}{The maximum speed of the particle will be, when $\phi=\cos^{-1}(\mu\cot\alpha)$}
\end{qoptions}
\end{bankq}
\begin{bankpassage}[id=STAR01-P1-PHY-S2-G7-8, source={Star Review (Other) - Test 01 (12 Sep 2022)}, questions={7-8}]
A small block of mass $M$, starts sliding down on a bigger block of mass $2M$, having a circular track of radius $R$. Let $N_1$ and $N_2$ be the normal reactions on the bigger block due to the wall and the floor. Do not consider the motion of the small block after $\theta=90^\circ$. (All surfaces are frictionless)
\bfig{142.8}{STAR01_p1_phy_s2_para7_8_1.png}
\end{bankpassage}
\begin{bankq}[id=STAR01-P1-PHY-S2-Q07, source={Star Review (Other) - Test 01 (12 Sep 2022)}, paper=1, qno=7, type=single, subtopic={Constraints}, answer={B}, answer_src=key, difficulty=6, diff_parts={K1 S2 I2 C2 F0}, diff_note={concepts: small block sliding on a larger block with a circular track | steps: speed of the small block at the angle; normal force; wall reaction from the horizontal force balance; evaluate at thirty degrees}, group=STAR01-P1-PHY-S2-G7-8]
The value of $N_1$ at the instant when $\theta=30^\circ$
\begin{qoptions}
\opt{A}{$\dfrac{\sqrt3Mg}4$}
\opt{B}{$\dfrac{3\sqrt3Mg}4$}
\opt{C}{$\dfrac{\sqrt3Mg}2$}
\opt{D}{$\dfrac{3\sqrt3Mg}2$}
\end{qoptions}
\end{bankq}
\begin{bankq}[id=STAR01-P1-PHY-S2-Q08, source={Star Review (Other) - Test 01 (12 Sep 2022)}, paper=1, qno=8, type=single, subtopic={Constraints}, answer={C}, answer_src=key, difficulty=7, diff_parts={K1 S3 I2 C2 F0}, diff_note={concepts: small block sliding on a larger block with a circular track | steps: speed at the angle; normal force; vertical balance of the large block; set the floor reaction; solve for the angle}, group=STAR01-P1-PHY-S2-G7-8]
The value of $\tan\theta$ at the instant when $N_2=3\,Mg$
\begin{qoptions}
\opt{A}{$\dfrac1{\sqrt2}$}
\opt{B}{$\sqrt2$}
\opt{C}{$\dfrac12$}
\opt{D}{$2$}
\end{qoptions}
\end{bankq}
\begin{bankpassage}[id=STAR01-P1-PHY-S2-G9-10, source={Star Review (Other) - Test 01 (12 Sep 2022)}, questions={9-10}]
A ball of mass $M$ attached to a string of length $\ell$ moves in a circle in a vertical plane as shown. At the top of the circular path, the tension is twice the weight of the ball. At the bottom, the ball just clear the ground.
\bfig{134.9}{STAR01_p1_phy_s2_para9_10_1.png}
\end{bankpassage}
\begin{bankq}[id=STAR01-P1-PHY-S2-Q09, source={Star Review (Other) - Test 01 (12 Sep 2022)}, paper=1, qno=9, type=single, subtopic={Circular dynamics}, answer={D}, answer_src=key, difficulty=4, diff_parts={K1 S2 I1 C1 F0}, diff_note={concepts: vertical circle with a given tension at the top | steps: top speed from the tension; energy conservation; lowest speed}, group=STAR01-P1-PHY-S2-G9-10]
Determine the speed of the ball at the lower most point.
\begin{qoptions}
\opt{A}{$\sqrt{3g\ell}$}
\opt{B}{$\sqrt{5g\ell}$}
\opt{C}{$\sqrt{6g\ell}$}
\opt{D}{$\sqrt{7g\ell}$}
\end{qoptions}
\end{bankq}
\begin{bankq}[id=STAR01-P1-PHY-S2-Q10, source={Star Review (Other) - Test 01 (12 Sep 2022)}, paper=1, qno=10, type=single, subtopic={Circular dynamics}, answer={A}, answer_src=key, difficulty=4, diff_parts={K1 S2 I1 C1 F0}, diff_note={concepts: projectile after the string is cut at the top | steps: top speed; height of the point; time of fall; horizontal distance}, group=STAR01-P1-PHY-S2-G9-10]
If the string is cut when the ball is at the top at $t=t_0$, then the horizontal distance the ball travels before hitting the ground and after $t=t_0$ is :-
\begin{qoptions}
\opt{A}{$2\sqrt3\ell$}
\opt{B}{$4\sqrt3\ell$}
\opt{C}{$3\sqrt3\ell$}
\opt{D}{$5\sqrt3\ell$}
\end{qoptions}
\end{bankq}
\begin{bankpassage}[id=STAR01-P1-PHY-S3-G1-2, source={Star Review (Other) - Test 01 (12 Sep 2022)}, questions={1-2}]
A man of mass 50 kg stands on a horizontal weighing machine of negligible mass, attached to massless platform that slide down a $37^\circ$ rough incline plane. The reading of weighing machine is 48 kg (take : $g=10$ m/s$^2$)
\bfig{92.6}{STAR01_p1_phy_s3_para1_2_1.png}
\end{bankpassage}
\begin{bankq}[id=STAR01-P1-PHY-S3-Q01, source={Star Review (Other) - Test 01 (12 Sep 2022)}, paper=1, qno=1, type=numerical, subtopic={Newton's laws}, answer={0.40}, answer_src=key, difficulty=5, diff_parts={K1 S2 I1 C1 F1}, diff_note={concepts: apparent weight of a man on a sliding platform | steps: normal force from the reading; vertical equation; vertical acceleration}, group=STAR01-P1-PHY-S3-G1-2]
What is the vertical acceleration (in m/s$^2$) of man ?
\end{bankq}
\begin{bankq}[id=STAR01-P1-PHY-S3-Q02, source={Star Review (Other) - Test 01 (12 Sep 2022)}, paper=1, qno=2, type=numerical, subtopic={Newton's laws}, answer={0.665}, answer_src=key, difficulty=5, diff_parts={K1 S2 I1 C1 F1}, diff_note={concepts: acceleration along the incline from the vertical component | steps: vertical acceleration of the platform; relation with the slope angle; acceleration along the slope}, group=STAR01-P1-PHY-S3-G1-2]
What is the acceleration (in m/s$^2$) of man down the slope $a_s$ ?
\end{bankq}
\begin{bankpassage}[id=STAR01-P1-PHY-S3-G3-4, source={Star Review (Other) - Test 01 (12 Sep 2022)}, questions={3-4}]
Masses of blocks A, B and C are 7kg, 4kg and 1 kg respectively. String and pulleys are ideal. All surfaces are smooth.
\bfig{116.5}{STAR01_p1_phy_s3_para3_4_1.png}
\end{bankpassage}
\begin{bankq}[id=STAR01-P1-PHY-S3-Q03, source={Star Review (Other) - Test 01 (12 Sep 2022)}, paper=1, qno=3, type=numerical, subtopic={Pulley systems}, answer={4.00}, answer_src=solved, working={C stays against A so $a_C^x=a_A$. With $T$ the string tension: A and C get $2T$ in total, $a_A=\dfrac T4$; B gets $2T$ opposite, $a_B=\dfrac T2$; C falls with $a^y=10-T$. String length constraint $2a_B-2a_A+a^y=0$ gives $T=4$ N}, difficulty=8, diff_parts={K1 S3 I2 C2 F1}, diff_note={concepts: blocks with a string over pulleys on moving blocks | steps: forces on each block; constraint relation of the string; contact of the small block; equations of motion; solve the tension}, group=STAR01-P1-PHY-S3-G3-4]
Find the tension in the string (in N).
\end{bankq}
\begin{bankq}[id=STAR01-P1-PHY-S3-Q04, source={Star Review (Other) - Test 01 (12 Sep 2022)}, paper=1, qno=4, type=numerical, subtopic={Pulley systems}, answer={37.00}, answer_src=key, difficulty=8, diff_parts={K1 S3 I2 C2 F1}, diff_note={concepts: blocks with a string over pulleys on moving blocks | steps: forces on each block; constraint relation of the string; contact of the small block; equations of motion; tension; acceleration components}, group=STAR01-P1-PHY-S3-G3-4]
The acceleration of block `C' is $\sqrt N$ m/s$^2$. Find `N'.
\end{bankq}
\begin{bankpassage}[id=STAR01-P1-PHY-S3-G5-6, source={Star Review (Other) - Test 01 (12 Sep 2022)}, questions={5-6}]
A sufficiently long plank of mass 4 kg is kept on a smooth horizontal ground. A block of mass 2kg kept over the plank and is being acted upto by a time varying horizontal force $F=0.5\,t$, where $F$ is in newton and $t$ is in seconds as shown in figure. The coefficient of friction between the plank and the block is `$\mu$'. At $t=12$ sec relative motion between the plank and the block just starts.
\bfig{95.0}{STAR01_p1_phy_s3_para5_6_1.png}
\end{bankpassage}
\begin{bankq}[id=STAR01-P1-PHY-S3-Q05, source={Star Review (Other) - Test 01 (12 Sep 2022)}, paper=1, qno=5, type=numerical, subtopic={Friction}, answer={0.2}, answer_src=key, difficulty=5, diff_parts={K1 S2 I1 C1 F1}, diff_note={concepts: plank and block with a time varying force | steps: limiting friction at the onset of sliding; acceleration of the plank at that time; coefficient}, group=STAR01-P1-PHY-S3-G5-6]
Find the coefficient of friction $\mu$.
\end{bankq}
\begin{bankq}[id=STAR01-P1-PHY-S3-Q06, source={Star Review (Other) - Test 01 (12 Sep 2022)}, paper=1, qno=6, type=numerical, subtopic={Friction}, answer={0.6}, answer_src=key, difficulty=6, diff_parts={K1 S2 I2 C1 F1}, diff_note={concepts: average acceleration of a plank over an interval | steps: motion before slipping; motion after slipping; speed at fifteen seconds; average acceleration}, group=STAR01-P1-PHY-S3-G5-6]
Find the average acceleration of the plank in the time interval 0 to 15 s in the figure (in m/s$^2$).
\end{bankq}
\begin{bankq}[id=STAR01-P1-PHY-S4-Q09, source={Star Review (Other) - Test 01 (12 Sep 2022)}, paper=1, qno=9, type=numerical, subtopic={Constraints}, answer={0}, answer_src=key, difficulty=5, diff_parts={K1 S2 I1 C1 F1}, diff_note={concepts: tension in a rope sliding on a smooth surface | steps: acceleration of the rope; mass below the point; tension}]
As shown in the arrangement a uniform heavy rope AB of mass 10 kg slides down smooth surface. Find the tension (in N) at $\dfrac34$th length of rope from point A.
\bfig{145.2}{STAR01_p1_phy_s4_q09_1.png}
\end{bankq}
\begin{bankq}[id=STAR01-P2-PHY-S1-Q09, source={Star Review (Other) - Test 01 (12 Sep 2022)}, paper=2, qno=9, type=single, subtopic={Circular dynamics}, answer={A}, answer_src=key, difficulty=7, diff_parts={K1 S3 I2 C2 F0}, diff_note={concepts: force on a particle at the minimum speed on a rod | steps: condition for the vertical circle; speed at the position; radial force; tangential force; net}]
A particle is attached at one end of a massless rod whose other end is fixed at $O$ as shown in figure. A particle is given minimum velocity at lower most point to complete vertical circular motion about $O$. Find net force on particle when it is at position $P$. Length of rod is $L$.
\bfig{76.8}{STAR01_p2_phy_s1_q09_1.png}
\begin{qoptions}
\opt{A}{$\dfrac{mg}2\sqrt{39}$}
\opt{B}{$\dfrac{mg}2\sqrt{33}$}
\opt{C}{$\dfrac{mg}2\sqrt{47}$}
\opt{D}{$\dfrac{2mg}3\sqrt{55}$}
\end{qoptions}
\end{bankq}
\begin{bankq}[id=STAR01-P2-PHY-S1-Q13, source={Star Review (Other) - Test 01 (12 Sep 2022)}, paper=2, qno=13, type=single, subtopic={Circular dynamics}, answer={C}, answer_src=key, difficulty=4, diff_parts={K1 S2 I1 C1 F0}, diff_note={concepts: minimum speed on a death well | steps: friction balance weight; normal force from the circular motion; speed}]
A rider wants to ride a motorcycle on the vertical surface of a large cylindrical wooden wall on commonly known as death well. The radius of well is 20 m and coefficient of friction between the tyres of motorcycle and wall is 0.5. The minimum speed motorcyclist must have in order to prevent slipping should be :
\begin{qoptions}
\opt{A}{30 m/s}
\opt{B}{25 m/s}
\opt{C}{20 m/s}
\opt{D}{15 m/s}
\end{qoptions}
\end{bankq}
\begin{bankq}[id=STAR01-P2-PHY-S1-Q16, source={Star Review (Other) - Test 01 (12 Sep 2022)}, paper=2, qno=16, type=single, subtopic={Friction}, answer={D}, answer_src=key, difficulty=4, diff_parts={K1 S2 I1 C1 F0}, diff_note={concepts: relative acceleration of two blocks with friction | steps: check for slipping; acceleration of each; relative acceleration}]
Find acceleration of lower block w.r.t. upper block, if both blocks were at rest initially.
\bfig{147.5}{STAR01_p2_phy_s1_q16_1.png}
\begin{qoptions}
\opt{A}{1 m/s$^2$}
\opt{B}{0.5 m/s$^2$}
\opt{C}{0.25 m/s$^2$}
\opt{D}{0 m/s$^2$}
\end{qoptions}
\end{bankq}
\begin{bankq}[id=STAR01-P2-PHY-S1-Q17, source={Star Review (Other) - Test 01 (12 Sep 2022)}, paper=2, qno=17, type=single, subtopic={Pulley systems}, answer={B}, answer_src=key, difficulty=4, diff_parts={K1 S2 I1 C1 F0}, diff_note={concepts: initial acceleration in a pulley system | steps: constraint between the blocks; equations of motion; solve}]
System is released at $t=0$. Find acceleration of block of mass $m$ at $t=0$.
\bfig{63.5}{STAR01_p2_phy_s1_q17_1.png}
\begin{qoptions}
\opt{A}{$\dfrac{3g}2$}
\opt{B}{$\dfrac{2g}3$}
\opt{C}{$\dfrac{8g}5$}
\opt{D}{$\dfrac{9g}4$}
\end{qoptions}
\end{bankq}
\begin{bankq}[id=STAR01-P2-PHY-S1-Q18, source={Star Review (Other) - Test 01 (12 Sep 2022)}, paper=2, qno=18, type=single, subtopic={Friction}, answer={B}, answer_src=key, difficulty=4, diff_parts={K1 S2 I1 C1 F0}, diff_note={concepts: minimum force to pull a block | steps: force at an angle; normal force; condition; minimise}]
A block of mass 4 kg is placed on a rough horizontal surface and coefficient of friction between block and surface is 0.75. Minimum force required to pull the block on horizontal surface is :
\begin{qoptions}
\opt{A}{15 N}
\opt{B}{24 N}
\opt{C}{30 N}
\opt{D}{40 N}
\end{qoptions}
\end{bankq}
\begin{bankq}[id=STAR01-P2-PHY-S1-Q19, source={Star Review (Other) - Test 01 (12 Sep 2022)}, paper=2, qno=19, type=single, subtopic={Friction}, answer={B}, answer_src=key, difficulty=4, diff_parts={K1 S2 I1 C1 F0}, diff_note={concepts: range of the hanging mass for equilibrium | steps: forces along the incline; limiting friction in each direction; maximum and minimum}]
A block of mass 10 kg is placed on rough inclined having angle of inclination $37^\circ$ \& coefficient of friction $\mu=\dfrac12$. This 10 kg block is connected with light string with an another mass $M$ as shown in figure. Initially system is at rest. Calculate maximum and minimum value of $M$ so system does not move.
\bfig{92.3}{STAR01_p2_phy_s1_q19_1.png}
\begin{qoptions}
\opt{A}{14 kg and 2 kg}
\opt{B}{10 kg and 2 kg}
\opt{C}{15 kg and 3 kg}
\opt{D}{9 kg and 1 kg}
\end{qoptions}
\end{bankq}
\begin{bankq}[id=STAR01-P2-PHY-S2-Q04, source={Star Review (Other) - Test 01 (12 Sep 2022)}, paper=2, qno=4, type=numerical, subtopic={Friction}, answer={6}, answer_src=key, difficulty=5, diff_parts={K1 S2 I1 C1 F1}, diff_note={concepts: force at which a block slides over another | steps: maximum acceleration of the upper block; pull of the hanging mass; solve}]
$m_1=4$ kg, $m_2=2$ kg, $\mu=0.4$ \& $g=10$ m/s$^2$. Find $M$ (kg) for which $m_2$ will start sliding over $m_1$. (Ground is frictionless)
\bfig{147.5}{STAR01_p2_phy_s2_q04_1.png}
\end{bankq}
\begin{bankq}[id=STAR01-P2-PHY-S2-Q05, source={Star Review (Other) - Test 01 (12 Sep 2022)}, paper=2, qno=5, type=numerical, subtopic={Friction}, answer={12}, answer_src=solved, working={Ground friction limit $0.6\times200=120$ N. The upper block slips when the common acceleration is $\mu g=6$: $F-120=20\times6$, $F=240$ N, $20t=240$, so $t=12$ s}, difficulty=6, diff_parts={K1 S2 I2 C1 F1}, diff_note={concepts: slipping of stacked blocks under a growing force | steps: limiting friction with the ground; maximum acceleration of the upper block; force at slipping; time}]
In the given figure, a force of magnitude $20\,t$ is applied on lower block where $t$ is time in sec. Coefficient of static friction between contact surface is 0.6. For what value of `$t$' upper block begins to slip relative to lower block ?
\bfig{150.1}{STAR01_p2_phy_s2_q05_1.png}
\end{bankq}
\end{document}
